\section{A separate Grothendieck dictionary: objects, operations, and proof transport}
\label{sec:grothendieck-dictionary}

This part restates the checked deductions in the requested geometric language. A locator identifies the framework in which a proof lives; it is not a claim that Grothendieck stated a modern complexity result. The original manuscript images are not in the attachment. Cote and scan references below therefore refer to the supplied transcriptions, and the mathematical corrections do not certify an unseen handwritten reading.

\subsection{Exact framework locators and scope}

\paragraph{D-S3: group schemes and weights.}
For an affine group scheme, its representations are comodules over its Hopf algebra. The relevant reference is SGA~3, Expos\'e~I, Proposition~4.7.2. For a diagonalizable group $D(M)$, such comodules are $M$-graded modules, by Proposition~4.7.3; the following remark gives the corresponding graded-algebra description of affine schemes with action. Proposition~4.7.4 gives equivariant splittings when the underlying sequence splits. These statements, inspected in the original Expos\'e~I PDF, give the integral meaning of every torus weight and cocharacter below. For stabilizers, SGA~3, Expos\'e~VI$_B$, Proposition~4.1 gives upper semicontinuity of group-fiber dimension and explicitly cites EGA~IV$_3$, Theorem~13.1.3. The latter locator is corroborated through that primary SGA text; its entire EGA proof has not been reread here.

\paragraph{D-S6: $K_0$, exterior powers, Adams operations, and filtration.}
SGA~6, Expos\'e~V, Definition~2.1 and Example~2.2(b), construct the pre-$\lambda$ structure of locally projective modules over a ringed topos; Lemma~2.2.1 uses a filtration on exterior powers of an exact sequence. Definition~2.4 and equation~(2.4.2) specify the universal product and composition identities. Expos\'e~VI, section~3.2 gives the special-$\lambda$-ring assertion for the $K'$ of a quasi-compact scheme. Expos\'e~V, Proposition~3.2 defines $\gamma$ from $\lambda$; section~3.10, equation~(3.10.1), defines the augmentation filtration. Section~7.1, equation~(7.1.1), and Proposition~7.5 give Adams operations and their ring identities. The electronic transcriptions of the cited Expos\'e~V passages were inspected; the exact VI~3.2 statement was corroborated in the primary-text transcription during independent review. We do not describe an integral character identity as an integral semisimple decomposition.

For split reductive group schemes over a principal ideal domain, in particular over $\mathbb Z$, a useful complementary reference is Serre (1968), Theorem~5, section~3.7, and the explicit $\mathrm{GL}_n$ character ring in section~3.8. The base-change statements in that theorem concern the representation Grothendieck group, not preservation of a characteristic-zero basis of simple objects in every characteristic. Nor do they automatically answer the separate finite-symmetric-group lattice question recorded in the supplied study of cote~133.


\paragraph{D-Tor: derived restriction and its obstruction groups.}
SGA~6, Expos\'e~III, Definition~1.5 defines Tor independence of two
schemes over a base by vanishing of their positive Tor sheaves; in
the affine case it is exactly the assertion that derived tensor
product agrees with ordinary tensor product. Notation~1.6 specifies
the relative derived tensor product. Expos\'e~I, Proposition~5.1
and Definition~5.2 give the flat-amplitude criterion, and paragraph~5.4
identifies it with ordinary finite flat-resolution length for a
module. These exact passages were read in the electronic French
primary-text transcription. III~1.5 refers to EGA~III, section~6.5,
for local hyper-Tor; that reference is corroborated here through
SGA~6, without a new whole-section reading of EGA~III.

\paragraph{D-Proper: closed images of projective incidence schemes.}
EGA~II, Definition~5.4.1, printed p.~100, defines properness
using separated finite type and universal closedness.
Theorem~5.5.3(i), printed p.~104, states that a projective
morphism is proper. The corresponding French primary-text
passages, including the displayed statements and the proof
following the latter, were read through the indexed original
PDF text. The complete PDF did not load in the web reader.
These locators justify the closed incidence images below;
they do not supply a computational bound for their equations.

\paragraph{D-FGA: images, quotients, and projective parameter spaces.}
For a homogeneous orbit map over $\mathbb Z$, take its schematic image by the kernel of the coordinate homomorphism, as in G-04. Degreewise kernels and images are given by finite matrices and their minors. EGA~II, sections~2.4--2.5, in particular Proposition~2.5.2, gives the projective scheme and associated sheaf of a graded module. The electronic EGA translation was inspected for this locator. FGA, Bourbaki Expos\'e~221, Theorem~3.1 gives the Quot parameter space, Theorem~3.2 its fixed-Hilbert-polynomial projective form, and Proposition~3.8 its Grassmannian realization in the required degree range. The original FGA scan was inspected. These results place the families in parameter spaces; they supply no polynomial bound for finding the relevant minors or their nonvanishing loci.

\paragraph{D-ET: traces, weights, and their arithmetic size.}
SGA~4$\frac12$, \emph{Rapport sur la formule des traces}, Theorems~3.1--3.2 and Proposition~3.3, gives the determinant expression for the $L$-function, the trace formula, and the logarithmic-derivative identity. SGA~5, Expos\'e~X, Theorem~7.1 and Corollary~7.12, formula~(7.13), gives the Euler--Poincar\'e conductor accounting. Deligne's Weil~II, Theorem~3.3.1 and section~3.3.10, supplies the compact-support weight estimate used in P-06; Theorem~3.5.3 gives the equidistribution setting. The Weil~II locators were checked against Katz and Sarnak (1999), sections~9.0.13--9.0.16 and Theorem~9.2.6. The original large Weil~II PDF was not successfully opened here. This is a statement-and-dependency check, not an independent verification of the complete proof of Weil~II.

\subsection{Rectangular characters and orbit-closure multiplicities}

\paragraph{R-01. Tensor pullback of a Schur operation.}
For locally free $E,F$ of ranks $a,n$, the tensor morphism $\mathrm{GL}(E)\times\mathrm{GL}(F)\to\mathrm{GL}(E\otimes F)$ is defined over $\mathbb Z$. The class $s_{(d^n)}[E\otimes F]$ restricts, after passage to characteristic zero, to the sum of $[S_\mu E\otimes S_\nu F]$ with coefficients $g((d^n),\mu,\nu)$. The Haar projector is the complex realization of extraction of that character component. Its finite CUE dimension is a rank of a Schur construction; its normalized limiting constant need not be an integer rank. This is precisely the proof of R-01 transported through D-S3 and D-S6. Cote~133, scan~028, lines~6--17, provides the tensor-power and symmetric-operation language, while scan~019 distinguishes rank from the full character class.

\paragraph{R-02. Two different meanings of positivity.}
The tableau expansion is an integral character identity in the universal special $\lambda$-ring, with the Weyl alternating factor extracting a highest weight. Over $\mathbb C$, the invariant summand of $S_\lambda(E\otimes F)\otimes\det(E)^{-d}\otimes\det(F)^{-d}$ is an actual vector space. A Hermitian metric identifies its projector trace with a sum of squared absolute values. The former is an integral identity; the latter uses complex conjugation and compact integration. D-S6 and cote~133, scan~028, lines~19--24 (the direct-summand question), are the exact conceptual distinction. No cancellation-free combinatorial basis is inferred from the analytic squares.

\paragraph{R-03. Central quotient, disconnected stabilizer, and rank.}
The tensor morphism has the central subgroup $t\mapsto(tI,t^{-1}I)$ in its kernel, a statement of group schemes over $\mathbb Z$. A compact block-Levi branching rule is therefore a different restriction functor. The transpose component of the determinant stabilizer is additional group-scheme data beyond the connected tensor subgroup. The $S_4$ counterexamples compare distinct invariant functors and survive as finite character computations. D-S3 places these groups and actions; cote~133, scan~019, lines~45--56, is the corresponding warning that augmentation loses equivariant information. The negative Weingarten coefficient is a complex integration fact, not a negative rank of a locally free sheaf.

\paragraph{R-04. Plethysm as an operation on representation $K_0$.}
The cycle-index formula expresses $\sigma^d(\sigma^n[E])$ using the Adams operations $\psi^j$. Its displayed rational coefficients are interpreted in $K_0\otimes\mathbb Q$ or the universal torsion-free character ring; the resulting symmetric-power operation is integral. D-S6, V~7.1 and V~7.5, makes their algebraic status precise. They are ring operations; $g\mapsto g^j$ need not be a group homomorphism, and $\psi^j$ need not preserve the cone of honest representations. The $U(2)$ mixed-moment counterexample therefore translates to the statement that separate character distributions do not determine the joint operation. The wreath-subgroup invariants give the characteristic-zero tensor-functor realization. Cote~133, scan~028, lines~12--17 and~26--35, places tensor powers and induction products, with each $\rho_i$ in its own homogeneous degree.

\paragraph{R-05. Image modules, not merely stabilizer invariants.}
The intrinsic object is the image of the degree-$d$ orbit-coefficient map. Its dual is the span of the orbit of $f^{\otimes d}$. Over $\mathbb C$, complete reducibility identifies the multiplicity in each Schur summand with the rank of its multiplicity matrix. The Haar covariance and its reduced Gram matrices have exactly those images; their principal exterior coefficients are squared Pl\"ucker norms. This transports the entire rank proof to determinantal rank strata and Grassmannians in D-FGA. The Hermitian form is an auxiliary real structure and is not asserted to define a motive. Cote~133, scan~028's question about a direct summand is answered here only after passing to characteristic zero.

\paragraph{R-06. Integral matrices and special fibers.}
The polynomial orbit map is defined over $\mathbb Z$ after the specified symmetrizer scalings. Its finite coefficient matrix is an integral map of free modules. The Gaussian covariance is another complex Gram form on its image, while the integral character-sum matrix computes the required Schur-isotypic rank after scalar extension. The Schur projectors require the stated factorial denominators. G-04 specifies the open localization on which the ranks and kernels specialize correctly. D-FGA records the associated rank and Quot strata; D-S6 records the tensor and Schur operations. This translation preserves the exponential matrix sizes, so changing language does not remove the computational gate.

\paragraph{R-07. Binary restriction and highest weights.}
The determinant-pencil substitution is a morphism of coefficient schemes, defined after clearing the displayed binomial denominators. The sections $Q_j$ are highest-weight polynomial functions, and independence is proved by their restrictions to the chosen diagonal-pencil chart. Over $\mathbb C$ every binary form is a product of linear forms, so the binary coefficient subscheme lies in the determinant orbit closure. Highest-weight sections of length at most two therefore cannot belong to its ideal. This is the same proof expressed as ideal restriction in D-FGA and weights in D-S3. Cote~133, scan~028 supplies the polynomial-functor viewpoint; it does not supply a new multiplicity formula by rank alone.

\paragraph{R-08. The quotient class in equivariant $K_0$.}
For the polynomial ring $S=\operatorname{Sym}(W)$ and orbit quotient $A=S/I$, an equivariant free resolution yields $[A]=\sum_i(-1)^i[S\otimes T_i]$ in the appropriate graded Grothendieck group. Symmetric powers produce the Molien denominator and the alternating resolution produces the numerator. Exactness, not pointwise sign, proves the formula. D-S6, V~2.2(b), and D-FGA place the additivity and coherent quotient. Cote~133, scan~019, lines~9--20, contains the exterior/symmetric generating-series relation. Ignoring the numerator forgets the defining ideal and is not an operation allowed by $K_0$ additivity.

\paragraph{R-09. A four-factor invariant and a finite-rank cutoff.}
The squared rectangular multiplicity sum is the determinant-isotypic component of $\operatorname{Sym}(E_1\otimes E_2\otimes E_3\otimes E_4)$ in degree $nd$, obtained by composing two tensor restriction functors. The $U(N)$ cutoff means that Schur functors with more than $N$ rows vanish on a rank-$N$ bundle. The stable ratios denominator is therefore an identity of symmetric operations on four tensor factors. D-S3 and D-S6 carry the proof without an arithmetic purity hypothesis. Cote~133, scan~028, is the corresponding tensor-power locus. Nonnegative character norms are exact, but no effective algorithm for the remaining coefficient extraction has been added by the translation.

\paragraph{R-10. The quadratic case and the component group.}
Over $\mathbb C$, the $2\times2$ determinant is a nondegenerate quadratic form in four variables. Its orbit closure is the whole quadratic-form coefficient space. The connected group is the tensor image of $\mathrm{SL}_2\times\mathrm{SL}_2$; the full quadratic stabilizer includes the other orthogonal component. Accordingly the all-odd four-row part present in the ordinary rectangular Kronecker coefficient is removed from the symmetric coefficient and the coordinate ring. D-S3 records the group and component distinction, while D-S6 identifies the even-part Schur decomposition. Cote~133, scan~019's augmentation does not retain this distinction.

\paragraph{R-11. Low-degree surfaces and restricted ideals.}
The linear and quadratic determinant constructions are coefficient-scheme morphisms. The cubic statement uses a dense open set of smooth cubic surfaces and Beauville's determinantal representation theorem; closedness extends the containment. Block padding is a morphism over the coefficient ring, and the highest-weight ideal inclusion is the same D-S3/D-FGA argument as G-08. It excludes four-row multiplicity obstructions for padded forms of degree at most three without asserting a global orbit containment for arbitrary many-variable cubics. Cote~133, scan~028 provides the Schur-functor interpretation of the resulting coordinate modules.

\paragraph{R-12. A Jacobian certificate and its characteristic-zero consequence.}
The permanent-pencil coefficient map is a morphism of affine schemes over $\mathbb Z$. A maximal Jacobian minor nonzero modulo the stated prime is a nonzero integer and therefore nonzero over $\mathbb Q$. In characteristic zero the differential-rank criterion gives dominance, so the four-variable projection of the permanent orbit closure fills the coefficient space in the certified sizes. D-FGA then translates the hypersurface quotient sequence of G-07 into the exact four-row multiplicity difference. The finite-field calculation is a certificate for an integer minor; it is not an appeal to equidistribution or an assertion of semisimplicity modulo that prime. The tensor-power language is cote~133, scan~028; the rank/character distinction remains scan~019.


\paragraph{R-13. A finite family of integral rank certificates and a torus quotient.}
For each stated size the Jacobian is a map of free modules on the affine pencil-parameter scheme. Its displayed minor defines a principal open rank locus over $\mathbb Z[1/101]$ and over $\mathbb Q$. The independent permanent identities check that these are the derivatives of the same morphism, rather than arbitrary matrices. For larger sizes the row-column torus acts in the fibers with the computed one-dimensional stabilizer. D-S3, VI$_B$~4.1 and the cited EGA dimension result, together with the elementary fiber-dimension calculation, give the upper bound. Highest-weight ideal restriction in D-FGA then yields the stated finite-range multiplicities. This extends the rank and Schur-functor dictionary of cote~133, scan~028; it does not turn a modular certificate into a characteristic-free semisimple decomposition.

\paragraph{R-14. The projective multiplication image and Hilbert growth.}
The map sending a line and a degree-$m$ form to their padded product is a morphism from $\mathbb P^3\times\mathbb P(\mathcal P_{m,4})$. Properness makes its image closed, and the unique linear factor of a general product makes the morphism generically injective. The resulting affine cone has dimension $c_m+3$. Its homogeneous ideal and Hilbert function are the D-FGA objects of EGA~II~2.4--2.5 and FGA Expos\'e~221, Theorems~3.1--3.2. The dimension comparison forces a larger Hilbert function in large degree, and D-S3 identifies the corresponding highest-weight multiplicities. The same proof therefore transports exactly; no effective degree threshold or new purity statement is inserted by this reformulation.

\paragraph{R-15. A normal direction to a family of pencil images.}
The cyclic pencil is a point of the affine parameter scheme, and $W_0$ is the image of its tangent map. Restriction to the explicitly defined lines $\Lambda_i\subset\mathbb P^3$ and the independent interval sections computes a lower bound for that image. The residue calculation proves that the section $T$ gives a nonzero class in the quotient by $W_0$. The two-chord family makes this class lie in the tangent image away from the special parameter, while an augmented determinant witnesses persistence of independence on a principal open. These are explicit maps of coefficient modules and open rank loci in D-FGA, defined over an integer coefficient ring after choosing the stated integral parameters. The generic characteristic-zero dimension and Hilbert argument then gives the four-row inequality for every equal size. D-S3 and D-S6 identify its weight and Schur class. Cote~133, scan~028 supplies the tensor-module setting; scan~019 explains why this is a multiplicity assertion rather than a numerical CUE dimension alone.

\subsection{Isolation as a torus action and a problem of coefficient modules}

\paragraph{I-01. Independent variable coordinates.}
The vertex-variable path polynomial is a section of the structure sheaf of an affine toric coordinate space. Its path monomials are distinct characters of its split torus. A Hoza weight is a cocharacter, and uniqueness of the minimum produces a one-dimensional initial character with a nonzero coefficient. D-S3 identifies the action and grading; I-10 supplies its flat hypersurface family over a field. Cote~133, scan~019, is only a finite-set analogy here, not a claim that path existence already computes a representation multiplicity.

\paragraph{I-02. Descent through a monomial quotient.}
The monomial map is a map of affine toric schemes whose character-lattice map is the integer matrix $A$. A cocharacter is induced by an integral weight on the full variable torus exactly when its degree vector lies in the image of $A^{\mathsf T}$. Homogeneity of the binomial kernel gives the grading on the image algebra; extending it to the full variable torus additionally requires the integral image-lattice condition in I-02. For example $z\mapsto x^2$ has zero kernel, but weight one on $z$ is not induced by an integer weight on $x$. The diamond is an explicit failure even of kernel homogeneity. This is D-S3, I~4.7.3, applied to the quotient algebra. Cote~133, scan~019 supplies no rule permitting independent path gradings to be identified with the quotient grading.

\paragraph{I-03. A complete coefficient flag.}
Filter the coefficient module by the span of monomials below each cocharacter weight. Basis isolation says that at each growth degree precisely one indexed coefficient survives, and that these surviving degrees are distinct. The dual formulation quantifies over every nonzero quotient functional. This is a statement about filtered locally free modules over a field and their graded pieces, not just the presence of one initial path. D-S3 gives the associated grading, and D-FGA places the open minor conditions on the relevant flag and Grassmann schemes. Cote~133, scan~028's direct-summand question motivates the distinction between a filtration and a chosen splitting.

\paragraph{I-04. The Pl\"ucker weight polytope.}
Take the coefficient row plane in a Grassmannian and its nonzero Pl\"ucker coordinates. Their torus weights are the sums of the corresponding monomial exponents. The vertex-fiber assertion is the symmetric-exchange proof stated in I-04; the probabilistic estimate counts walls crossed in one coordinate. The resulting open cones are normal cones refined by basis-weight inequalities. D-S3 identifies their cocharacters and D-FGA, Expos\'e~221, Proposition~3.8 places the Grassmannian. No finite-field purity theorem chooses an integer point in all these cones automatically.

\paragraph{I-05. Real Gram positivity and algebraic minors.}
With two independent dual coefficient matrices, Cauchy--Binet is an integral identity of sections of a determinant line. On the real locus where one is the conjugate transpose of the other, it becomes a sum of nonnegative squared Pl\"ucker coordinates. The counterexample shows that this determinant line does not retain the whole coefficient-image flag after colliding weights. D-FGA supplies the rank-stratum language, and D-S6 supplies the determinant as top exterior power. Cote~133, scans~019 and~028, respectively mark the distinction between numerical positivity and actual tensor-module information.

\paragraph{I-06. Secant rank of an evaluation functional.}
An evaluation at a point is a rank-one decomposable functional on the bounded-degree polynomial space. The moment matrix of the proposed nonlinear exponent functional has tensor-product rank. Expressing it as a fixed sum of evaluations would place it in a secant locus of bounded rank; the computed determinant excludes that locus below the stated exponential size. This is an algebraic rank statement over $K(t)$, hence belongs to the determinantal geometry in D-FGA. D-S3 also explains why a nonlinear exponent function is not a character-lattice homomorphism. No motive is needed for this counterexample.

\paragraph{I-07. Taking invariants changes the inserted data.}
The invariant functor on complex representations is realized by the Reynolds projector. An independent twirl sends each state insertion to its invariant scalar component. The two-insertion calculation resolves the commutant into the identity and transposition operations and exhibits a negative coefficient. Thus invariants are extracted in the representation category of D-S3/D-S6, with the full tensor operation retained. Cote~133, scan~028, describes why the symmetric-group action on tensor powers matters; scan~019 warns against replacing that action by dimension.

\paragraph{I-08. A finite quotient morphism, conditional on a hitting set.}
The trace-generating functions are regular invariant functions on the affine scheme of matrix tuples. A hitting set separating their coefficient polynomials gives a separating family in the quotient coordinate algebra; the cited invariant-theory theorem makes that algebra finite over the generated subalgebra. This is a finite morphism of affine schemes, not automatically a choice of algebraically independent parameters of minimal cardinality. D-FGA gives the general morphism and parameter-space setting. The unresolved explicitness condition is exactly the one retained in I-08 and GCT~V; the dictionary does not discharge it.

\paragraph{I-09. A tensor contraction computing a Hermitian norm.}
The matrix product computes a regular coefficient map. Tensoring with its conjugate and contracting equal monomial indices gives its squared norm; the torus integral extracts the zero character in that tensor product. D-S3 explains the character extraction. Over an integral algebraic base one must use independent dual variables, not treat complex conjugation as a regular morphism. The metric realization lives over $\mathbb R$ or $\mathbb C$. It is the same coefficient module as before, and access to its presenting matrices remains a white-box assumption.

\paragraph{I-10. Rees modules and specialization.}
The cocharacter gives an increasing filtration on coefficient spans. Its Rees module has the associated graded as its special fiber and the original module on the punctured line. The explicit hypersurface homogenization in I-10 is flat over a field by the torsion-free PID argument in its proof; general ideal degenerations require the entire weight-homogenized ideal. D-S3 identifies the grading and D-FGA identifies the flat family and its Hilbert/Quot points. The associated graded preserves total filtered rank, whereas adding colliding coefficients in a substituted polynomial can reduce their span, as I-05 proves. These are different maps.

\paragraph{I-11. Exterior powers do not remove equations.}
The exterior dimension and graph-of-linear-map count describe the sizes of coordinate and Grassmannian models. The two-term zero determinant is already an equation of the rank-deficient locus. Assigning different path weights to its cancelling terms changes that equation. D-S6 gives exterior powers and D-FGA gives the Grassmannian, but neither turns a signed determinant expansion into unsigned reachability. Cote~133, scan~028's direct-summand margin remains a relevant distinction between a functorial module and its termwise expansion.

\paragraph{I-12. Support is the nonvanishing locus of a regular function.}
A specified coefficient of an integer ROABP is a regular polynomial in its edge parameters. Its nonvanishing open subset is determined by a signed path sum, whereas ordinary reachability detects whether a parameterizing path set is empty. The GapL reduction is a computation of that regular function; the conditional complexity-class inclusion follows only under the specified logarithmic-space reduction. D-FGA locates its zero scheme. Cote~133, scan~019 supplies the nearby conceptual distinction between the cardinality of a finite set and a signed class; it states no complexity theorem.

\paragraph{I-13. Evaluation quadrature and the Segre variety.}
Coefficient tensors of width-one products form the affine cone over a Segre variety. Equality of the quadratic forms on all these tensors, with their complex conjugates, forces the moment-matrix identity proved in I-13. Its rank then bounds the number of evaluation vectors. This restates the entire proof as an assertion about a Segre cone, its dual evaluation vectors, and determinantal rank loci in D-FGA. The Hermitian positivity is not an integral $\lambda$-ring assertion. The much smaller hitting set in the same finding concerns only kernel avoidance, not equality of quadratic forms.


\paragraph{I-14. A Grassmannian open condition for several substitutions.}
The row plane of the coefficient matrix is a point of $\operatorname{Gr}(r,F^s)$. A fixed family of cocharacters gives the integral linear map $P_{\mathcal W}^{\mathsf T}$. The desired locus is the open set where its restriction to the tautological rank-$r$ bundle is injective. Its complementary equations are the displayed linear forms in Pl\"ucker coordinates. D-FGA, Expos\'e~221, Proposition~3.8 is the Grassmannian framework, and D-S3 identifies the cocharacters. The example proves that this open set depends on coordinate values and can be larger than the union of the individual basis-isolation loci.

\paragraph{I-15. A finite character-lattice obstruction.}
The bounded seed decoder produces a subspace of cocharacters of the split torus. Pigeonhole counting gives two distinct characters whose difference annihilates every cocharacter in that subspace. Their binomial is therefore a nonzero regular function in the kernel of every pullback. This is precisely the character/cocharacter pairing of D-S3, I~4.7.3, with the finite integer bounds retained. The extension to bounded-degree Boolean decoders is a dimension bound for their cocharacter span. It does not forbid more general maps whose spans are larger.

\paragraph{I-16. A finite coefficient-state scheme and base-field points.}
Over $\mathbb F_q$ the coefficient row space has $q^w$ rational points, and each letter acts on it by a linear morphism. The state automaton records these morphisms exactly. The separate evaluation lower bound counts the rational points avoiding products of prescribed coordinate divisors. Both constructions live in affine schemes over the finite field; D-S3 gives the coefficient representation, while D-FGA distinguishes the finite set of rational points from the whole scheme or its extension-field points. Purity does not change the explicit counting inequality.

\paragraph{I-17. A trivialization on an invertibility open set.}
The nonzero determinant polynomials define an open set on which the constant transition maps are bundle isomorphisms. The displayed gauge change trivializes those constant maps. The nested coefficient row subspaces can then grow at most $w-1$ times, and their chosen sections have the stated support. This is a concrete filtered-module argument in D-S3 and a rank-stratum argument in D-FGA. The known AGKS invertible-factor result is cited separately; the proof establishes the particular permutation-invariant sampler stated here.

\paragraph{I-18. A unit minor over a discrete valuation ring.}
Over $F[[\varepsilon]]$, a maximal minor defines an open rank locus precisely when it is a unit. A nonzero nonunit meets the generic fiber and misses the special fiber. The explicit matrix family separates these two cases, and the ideal of sampled scalar values shows why saturation in the fraction field does not recover an original evaluation. This is the elementary local-module form of the determinantal loci in D-FGA. The distinction from a flat Rees degeneration is essential: flatness of an ambient family does not imply surjectivity of this particular compressed module map.

\paragraph{I-19. Rank-factorization charts and the coefficient field.}
On the principal open where a chosen minor is invertible, a rank-$r$ bundle map factors through a rank-$r$ bundle. The adjugate formula clears its denominator and is an integral polynomial identity. Passing to the function field of the cut-variable scheme makes the intervening blocks disjoint; it does not remove the cut variables from the original base. Multilinearity proves coefficient-span preservation under independent local substitutions. D-FGA supplies these rank charts and D-S3 the tensor contractions. The unresolved descent from the cut-variable function field is explicitly retained.

\paragraph{I-20. Bounded losses in a sequence of coefficient subspaces.}
The product of a nonzero coefficient and the selected pivot minors defines a principal open on the cut-coordinate affine scheme. The finite grid meets this open by the degree bound. After specializing there, the coefficient subspaces have total dimension loss at most $\kappa$; hence their total growth is at most $w-1+\kappa$. This is an exact rank and filtration argument on the same D-FGA charts as I-19, followed by the D-S3 coordinate grading. It yields the stated sampler with its parameter dependence, without asserting that arbitrary singular sequences have bounded total loss.

\subsection{Purity, arithmetic sampling, and the limits of numerical shadows}

\paragraph{P-01. A Gram form on a finite space of sections.}
The coefficient support spans a vector space of functions. Their pairwise correlations are entries of a Gram matrix; its spectral lower bound is the elementary estimate in P-01. In geometric applications these functions may be sections or Frobenius traces, but no sheaf is manufactured by naming the Gram matrix. D-ET applies only after an actual trace-function realization exists. The distinction from rank augmentation in cote~133, scan~019, is that all cross terms have been retained and bounded.

\paragraph{P-02. A Kummer-character construction.}
The signs arise from the rank-one quadratic Kummer sheaf pulled back by $t\mapsto t+i$ on the line with its zeros removed. Products over distinct indices remain governed by the explicitly squarefree polynomial, so the curve character estimate controls the off-diagonal trace sums. The altered values at the excluded points contribute the stated elementary error. D-ET, together with the exact Alon--Goldreich--H\aa stad--Peralta locator, carries the proof. This yields the stated sparse multilinear hitting set; it does not create a polynomial bound on arbitrary coefficient supports.

\paragraph{P-03. The kernel of restriction.}
The sign-cube counterexample is a nonzero section whose restriction vanishes on the selected finite subscheme; the common-affine repair is annihilated by the explicit Vandermonde section. These are regular-function identities and survive independently of improved point-counting estimates. They belong to the affine-scheme language of D-FGA. Cote~133, scan~019 is an analogy for information lost by passing to a smaller numerical image, not a source of these new counterexamples.

\paragraph{P-04. Pullback to a curve and its zero divisor.}
For the curve morphism specified by the rational functions $r_i$, the section $f(r_1,\ldots,r_n)$ has polar divisor bounded by the displayed $D$. If it is nonzero, its zeros have total degree at most $\deg D$. The trace formula and curve RH then provide enough rational points. D-ET controls the count; the pullback's nonvanishing is a separate scheme-theoretic condition. This proof uses curves and sheaves actually specified in the statement, so a motivic interpretation is available for their cohomology, not for an unspecified ROABP.

\paragraph{P-05. A curve family can lie in a small polynomial kernel.}
Dimension comparison in the spaces $H^0(C,\mathcal O_C(aD))$ produces an algebraic relation between the two chosen functions. Multiplying the relations for the finite curve family gives its annihilating section and the explicit bivariate ROABP. D-FGA places the image and its ideal; D-ET cannot remove a function from that ideal by increasing the field size. The relation is proved directly with its degree and width bounds and is not a lower bound against every geometric sampler.

\paragraph{P-06. Purity with a complete conductor budget.}
The relevant class is $[\mathcal F_i\otimes\mathcal F_j^\vee]$ in the Grothendieck group of constructible Weil sheaves. The compactly supported trace gives its Gram entry. Geometric irreducibility and inequivalence determine $H_c^2$, affine vanishing removes $H_c^0$, and Weil~II bounds $H_c^1$. SGA~5 supplies the displayed Swan-conductor expression for its dimension. This is exactly D-ET, with every rank and conductor retained. A claim of polynomial explicitness additionally requires that these data and the enumeration of the points be polynomially bounded. Cote~133, scan~028 supplies a tensor-category analogy, but not the missing sheaf realization.

\paragraph{P-07. Sparse cocharacters over the integer lattice.}
The weight family is an explicit finite set in $\operatorname{Hom}(\mathbb G_m,\mathbb G_m^n)=\mathbb Z^n$. The collision polynomial over $\mathbb F_p$ is nonzero by its degree and coefficient bounds, so a seed outside the union of its root sets makes the support characters distinct. D-S3 translates the construction exactly; only the elementary finite-field root bound is used. Thus the result is independent of an additional purity or RH hypothesis.

\paragraph{P-08. Arithmetic covers with a specified nonvanishing encoding.}
The common number field determines a finite \emph{\'etale} cover of $\operatorname{Spec}\mathbb Z$ after inverting its discriminant. Frobenius conjugacy classes are the arithmetic data on that cover. Effective Chebotarev produces a small unramified prime only after the stated common map from primes to evaluations and its nonvanishing property are supplied. D-ET provides the parallel finite-field trace perspective; it is not a proof of that arithmetic encoding. The abelian case uses the stipulated Dirichlet GRH, while the general stated Hecke extension remains conditional on Masullo. No claim about arbitrary automorphic $\mathrm{GL}_n$ GRH is made.


\paragraph{P-09. Positive Hermitian forms on a Segre cone.}
The width-one coefficient tensors are decomposable sections, forming a Segre cone. Applying the upper sampled-energy bound to the conjugate evaluation tensor bounds each weighted evaluation norm, while applying the lower bound to the monomial sections bounds their total trace. This is the complete rank-one tensor proof of the cardinality inequality. D-S6 identifies tensor products and D-FGA the Segre coefficient scheme. The inequalities themselves use a complex Hermitian realization and are not inequalities in integral $K_0$.

\paragraph{P-10. Projective incidence and avoidance of a determinantal cone.}
The tensor-train class is the closed cone defined by the stated flattening minors. Its projectivization and evaluation hyperplane sections give a projective incidence correspondence. Proper projection makes the bad tuples a closed set; successive hyperplane sections show this set is proper at the stated cardinality. For unknown order one uses the finite union of the permuted cones. This is D-FGA's projective and determinantal geometry with D-S3's coefficient grading. Compactness of the complex unit sphere gives a positive Hermitian lower bound only after kernel avoidance is known. Neither the projective argument nor the $K_0$ class bounds the integer heights of an explicit good tuple.

\paragraph{P-11. Pullback along an explicit arithmetic parameterization.}
The presentation $\alpha_i=A_i(\theta)/b$ specifies a morphism from an affine parameter scheme to the variable space after inverting $b$. Pulling back a polynomial to $G_f(U,T)$ and using the stated degree grid gives the unconditional forward construction. Conversely Lagrange interpolation and a cyclotomic point give a finite \emph{\'etale} arithmetic realization after inverting the discriminant; minimal-polynomial degree prevents a nonzero rational pullback from vanishing there. D-FGA describes the morphism and its kernel, and D-ET places the finite cover. The proof does not require a Frobenius estimate: the missing hypothesis is the explicit common nonvanishing map itself, rather than a further trace bound.

\paragraph{P-Z01 and P-Z02. Realizations of zero statistics.}
These are analytic consequences of the precisely stated RH predicates and supported explicit-formula theorems. Their expression by unitary integrals does not make the zeta zeros the eigenvalues of a constructed geometric Frobenius. D-ET explains the function-field comparison only in the families with an established sheaf and monodromy group. There is no cote~133 scan containing a classical-RH-to-GCT implication; scans~019 and~028 supply representation operations, not that missing implication.

\subsection{The integral controls G-01--G-12}

\paragraph{G-01. Burnside classes, representation classes, and rank.}
The Burnside-ring counterexample is an equivariant finite-set computation. Linearization sends a $G$-set to a permutation module, but the two-subset operation does not become exterior power: the $C_2$ sign in the proof witnesses this failure. In D-S6 the correct relation is $x^2=2\lambda^2(x)+\psi^2(x)$, and only augmentation replaces $\psi^2(x)$ by the rank. This is the precise mathematical correction to cote~133, scan~019, line~53 and the supplied study's equivariant reading.

\paragraph{G-02. The universal terminal variable.}
Specializing one factor to the tensor unit in the universal product identity forces the top elementary symmetric function to remain present. SGA~6, V~2.4, equation~(2.4.2), explicitly includes that terminal variable. The proof in G-02 is its necessity in the universal character ring, so the correction to scan~006, lines~25--26, is mathematical and independent of paleography.

\paragraph{G-03. Equal finite filtration data, unequal representations.}
The two sums of line characters are honest locally free representations of $\mathbb G_{m,\mathbb Z}$, whose difference is the $(D+1)$st power of an augmentation-zero element. D-S6, V~3.10, places the difference in the required filtration step. The associated-bundle construction transports this identity to any scheme carrying the line bundle; it does not force the bundles to remain distinct on every base. Thus finite Chern or $\gamma$ data cannot recover all the character coefficients needed in scans~019 and~028.

\paragraph{G-04. The integral image and its good rank locus.}
The schematic orbit closure is defined over $\mathbb Z$ by a kernel, and its image algebra is torsion-free. In a fixed degree a nonzero maximal minor gives an explicit principal open on which the coefficient matrix is equivalent to a split rank matrix. This is a direct proof of base change for the displayed module, with D-FGA providing the image, Grassmannian and Quot framework. At the excluded primes the flat integral model's fiber may differ from the orbit closure of the reduced point. This is why integral $K_0$ identities from scan~028 do not imply unchanged characteristic-zero multiplicities in all fibers.

\paragraph{G-05. Plane curves and the three-row sector.}
The nonempty complement of the theta divisor in $\operatorname{Pic}^{g-1}(C)$ supplies Beauville's line bundle, hence a linear determinantal equation. Density of smooth curves gives the scheme-theoretic containment needed for every ternary coefficient section. Highest-weight sections of length at most three belong to that coefficient subalgebra, so its determinant ideal has none. D-FGA records the ideal restriction and parameter spaces; D-S3 records the weights. The numerical positive count is not extracted from Chern classes of the curve alone.

\paragraph{G-06. Stabilizer dimensions and four-row equations.}
The polynomial matrix-pencil map is invariant under the explicitly stated subgroup scheme $H$ of $\mathrm{GL}_n^2$. Its generic stabilizer is the scalar subgroup by the concrete matrix calculation. D-S3, VI$_B$~4.1 and its cited EGA theorem justify the dimension behavior, and the resulting proper image has a nonzero ideal. D-FGA and highest-weight restriction give an equation of exactly four rows, using G-05 to exclude shorter weights. This supplies no nonvanishing on a particular padded form until that extra statement is proved.

\paragraph{G-07. A Cartier divisor and a determinant twist.}
The reduced determinantal-quartic locus is a prime hypersurface. Its equation spans the character $\det^D$, so its ideal is the line-character twist of the ambient graded algebra shifted by $D$. The exact sequence in G-07 is the equivariant $K_0$ statement of a principal Cartier divisor. D-S6 gives additivity, D-FGA gives the projective hypersurface and quotient, and D-S3 identifies the character. It produces the shifted plethysm subtraction exactly; the corrected reduced degree is a literature issue, not a new transcription of cote~133.

\paragraph{G-08. A universal morphism of coefficient schemes.}
The two monomial trees and their cofactor determinant give a matrix over the universal integer coefficient ring. Homogenization and block padding therefore commute with any base change. Over $\mathbb C$ the containment of all low-variable restrictions yields an inclusion of highest-weight ideal spaces, hence the multiplicity inequality. This is D-FGA applied to an explicit scheme morphism, with the weights read through D-S3 and the symmetric powers through D-S6. Its complexity bound is part of the construction, not a consequence of representability alone.

\paragraph{G-09. A local padding coordinate and smaller universal matrices.}
Once the nonzero padding form is one coordinate, dehomogenization reduces the number of affine variables by one. The shared exponent chains are again defined over the universal coefficient ring; their exact vertex count gives the stated finite size. Changing an arbitrary padding form to that coordinate is local on its nonzero coefficient charts. D-FGA carries these local morphisms and their closure statement, and the D-S3 highest-weight argument yields the improved row exclusion. The known asymptotic uniform-representation theorem is cited separately; neither the chart construction nor its GCT consequence requires RH.


\subsection{Further rectangular, rank-profile, and arithmetic consequences}

\paragraph{R-16. Several normal directions on one deformation.}
The cyclic pencil lies on the coefficient affine scheme; disjoint row-pair modifications define one morphism from the affine line into that parameter scheme. The exact permanent identities produce polynomial sections of its differential-image sheaf away from zero. Their special-fiber classes are separated by characters of the finite diagonalizable group of $m$th roots of unity after adjoining $\zeta$ and inverting $m$. Dual Jacobi--Trudi is the Schur operation in D-S6, equivalently the exact cote~133 symmetric/exterior dictionary. A nonzero maximal minor makes the normal directions independent on a principal open. D-FGA and D-S3 therefore carry the dimension and padding proof; the characteristic-zero conclusion does not assert uniform good reduction at every prime.

\paragraph{G-10. Projective degree controls the first forced Hilbert gap.}
The pencil map induces a rational morphism of projective coefficient schemes. Pullback of general hyperplanes and the degree of their intersections bound the image degree, including components of the base locus. Chardin's bound then compares the graded dimensions with a polynomial subalgebra supplied by linear Noether normalization. These are D-FGA's projective, graded-module and Hilbert frameworks, with associated sheaves as in EGA~II, Proposition~2.5.2. The explicit bound is geometric effectivity; no small-size representation of every graded module is inferred.

\paragraph{G-11. Rectangular classes in a reductive quotient.}
Restriction from $\mathrm{GL}_4$ to $\mathrm{SL}_4$ sends precisely the determinant characters among polynomial irreducibles to the tensor unit. In D-S3 and D-S6, the degree-$d$ invariant module is therefore the rectangular multiplicity space. Explicit moment-map witnesses establish a nonempty stable open locus over $\mathbb C$; its quotient has dimension reduced by $15$, by the fiber-dimension statement already located in D-S3. The invariant rings are quotients because the group is linearly reductive in characteristic zero. The explicit Veronese is a standard graded coordinate algebra, so D-FGA applies again to its Proj, degree and Hilbert function. The analytic compact-group witness proves the stated characteristic-zero fact; it is not asserted to define a positive integral motive or to make reductive invariants exact over $\mathbb Z$.

\paragraph{G-12. A cocharacter degeneration and its invariant ideal.}
The homomorphism $t\mapsto\operatorname{diag}(t^3,t^{-1},t^{-1},t^{-1})$ is an integral cocharacter of $\mathrm{SL}_{4,\mathbb Z}$, as in D-S3. Its weight filtration extends the padded form over the affine line whenever $3(n-m)\geq m$. Strict inequality gives the zero special fiber; equality gives a ternary residual form with its padding factor. The ternary determinantal containment of G-05 then implies containment of invariant ideals in characteristic zero. In D-FGA this is the specialization of a morphism and the inclusion of its ideal of functions; in D-S6 it selects the determinant-character summands. It excludes precisely the specified rectangular sector.

\paragraph{I-21. Rank strata over local Artin bases.}
Every interval coefficient product is a map of free modules. The proof works over finite local Artin rings and requires a unit maximal minor, so the rank-two image and its complement split on the corresponding determinantal chart. Projection growth and its last-zero reset control the necessary coefficient supports. The two selector stages pull these finitely many rank conditions back to one nonempty open parameter locus. D-FGA supplies the determinantal and Grassmannian language, while D-S3 translates the coefficient supports to character gradings. This gives the explicit unknown-order width-three morphism; it does not assume that an arbitrary rank stratum over a nonreduced base splits without a unit minor.

\paragraph{P-12. A finite flat arithmetic evaluation and its trace dual.}
The algebra $\mathbb Z[T]/(T^N-2)$ is finite free of rank $N$. After inverting its discriminant it is finite \emph{\'etale}; over $\mathbb Q$ it is a field by Eisenstein. The integral coordinate polynomials $A_i$ define the evaluation morphism. Their bounded-degree pullback does not annihilate any nonzero member of the specified ROABP class, by I-21 and the explicit monomial encoding. The dual basis for the trace pairing extracts the transformed rational coefficient rows, which span the original coefficient-row space. This is an actual span-preserving arithmetic realization, not an injectivity assertion for differences of arbitrary members of the nonlinear width-three class. D-FGA describes the affine finite morphism and its kernel, D-S3 the coefficient lattice, and D-ET the finite \emph{\'etale} cover. Frobenius distribution and additional purity bounds are unnecessary for this construction.


\paragraph{R-17. Identical tangent-image modules and distinct isotypic sectors.}
At the cyclic pencil, the determinant and permanent differential maps have the same image up to signs of their columns. Their rank equality, together with the group-scheme fiber bound, gives the exact determinant-image dimension. D-FGA and D-S3 transport this computation. The two finite padded pairs then compare Hilbert polynomials, while the cocharacter specialization of G-12 removes the determinant-character summands. Their remaining positive differences are necessarily nonrectangular Schur summands in D-S6.

\paragraph{R-18. Two compatible normal families on the same scheme curve.}
Both local permanent identities hold on the same one-parameter deformation. The character decomposition of the special-fiber quotient now has two independent sections per permitted mode. The four-column coefficient determinant is the Schur operation $s_{(4^b)}$, interpreted in D-S6; its root-of-unity specialization is checked by a Weyl alternant. D-FGA turns their joint nonzero minor into an open rank condition, giving the improved padding range. The simultaneous nature of the deformation is necessary: ranks obtained at unrelated pencils are not added.

\paragraph{G-13. A weighted filtration inside an ordinary projective embedding.}
Generators of the reductive invariant quotient have specified positive degrees. Its cumulative graded module embeds in the ordinary polynomial-degree filtration of their affine coordinate embedding. Homogenizing that embedding gives an integral projective closure, to which the D-FGA Hilbert and degree comparison applies. Algebraically independent homogeneous sections provide the lower filtered dimension. This reduces the finite first-gap bound without asserting equality of the two filtrations or generation by a small Veronese degree.

\paragraph{I-22. Ordered coefficient modules and the radical quotient.}
The actual ROABP coefficients occupy an ordered product subspace, which need not equal the associative algebra generated by all layer coefficients. Passing to the semisimple quotient loses radical-valued sections even when their specified boundary functional is nonzero. In D-FGA these are different linear maps and different kernels; in representation $K_0$ the composition factors alone do not recover the off-diagonal extension data. The explicit triangular examples show the missing information before any trace or purity estimate can be applied.

\paragraph{I-23. Cocharacters preserving a Pl\"ucker subscheme.}
The Grassmann cone is cut out by the Pl\"ucker ideal. A diagonal cocharacter preserves that subscheme precisely when its weights satisfy the proved exchange equalities, which identify the integral torus weights induced from the original rank-$w$ bundle, up to a common character. D-S3 supplies the cocharacter lattice and D-FGA the Grassmannian. The binary classification and code-distance sum then quantify the cost of imposing this exact subscheme-preservation condition on every Hoza seed.

\paragraph{I-24. Constant-rank maps with repeated loss of remembered directions.}
The nested projections and their univariate perturbations define maps of free modules over an affine parameter scheme. Every interval lies on its specified unit-minor rank chart, including over the finite Artin quotients used in I-21. Nevertheless the selected boundary section is one monomial of exponential support. D-FGA's rank stratification records this phenomenon exactly; the ranks alone do not bound the length of a coefficient-support filtration.

\paragraph{I-25. Nested projective images on a Grassmannian.}
Invertible factors act regularly on the Grassmannian of rank-$r$ quotients or subspaces. The irreducible prefix image closures are nested because the constant factor is the identity; every strict inclusion raises dimension. D-FGA's Grassmannian chart gives dimension $r(w-r)$, and its projective geometry identifies scalar vanishing with containment in a hyperplane. Removing equal-image steps therefore gives the stated support existence bound. The invertible-factor hypothesis is what allows this proof to omit the reset behavior in I-24.

\paragraph{R-19. Compatible row and column image sections.}
Three-row motifs retain both supported row and supported column factorizations in one parameter family. Their four normal sections are separated by a six-column character calculation, whose Toeplitz determinant is the Schur operation $s_{(6^b)}$. D-S6 gives that operation, D-S3 the diagonalizable character splitting, and D-FGA the common nonzero determinantal minor. The resulting improved range is a statement about one actual pencil deformation; no addition of ranks from incompatible families is used.


\paragraph{G-14. Truncated coefficient jets and Schur-weight exclusion.}
The restriction modulo the square of the padding coordinate defines a linear
jet projection between affine parameter schemes. At a cyclic ternary pencil,
the differential is surjective: its two blocks are the cofactor-span map and
its multiplication by linear forms. The nonzero differential minor gives a
dominant morphism in characteristic zero (D-FGA). Every weight in the Schur
functor of shape $\lambda$ uses any fixed basis direction at most
$\lambda_1$ times, by its exterior-column construction in D-S6. After padding,
the low-weight part of an equation therefore factors through the first jet.
Its vanishing on the dominant determinant-jet image forces vanishing on all
padded forms. This is containment of the corresponding isotypic ideal
submodules, followed by exact passage to the characteristic-zero quotient;
it is not an assertion that the full padded varieties are contained.

\paragraph{I-26. Determinantal equations of an interpolation image.}
The maps $U$ and $V$ define bundle maps with $UV^T=0$. The determinant of
$U\operatorname{diag}(X)V^T$ is a nonzero section by a maximal-minor chart,
but its pullback to the centered interpolation parameter scheme vanishes
because that pulled-back bundle map factors through a rank-$k$ bundle.
The common center acts trivially. All exterior minors form a finite free
module on which rank-one additive translations act by commuting unipotent
matrices; the proof uses polynomial identities and remains valid in the
stated characteristics. This is the D-S3 additive group action and the
D-FGA determinantal rank locus. The surviving column-matroid bases give the
separate cocharacter nonvanishing certificate of D-S3. No purity assertion
turns the annihilated interpolation pullback into a nonzero section.


\paragraph{G-15. Harmonic covariants make the third-jet morphism dominant.}
The $\mathrm{SL}_2$ representation-ring decomposition of the endomorphism
bundle and the symmetric algebra of its adjoint bundle gives the covariants
$T_l$ in D-S3/D-S6. Their scalar coefficients in the adjugate are determined
by a finite trace pairing and its Jacobi matrix. The strictly positive
squared norm in the proof is an ordinary Hermitian linear-algebra fact,
not a weight-purity assertion. Linear multiples of these cofactor sections
surject onto the quotient by the fourth power of the padding-coordinate
ideal. Thus the differential minor gives dominance of the corresponding
third-jet morphism in D-FGA, while dimension excludes the next full jet.
The D-S6 exterior-column weight bound transfers dominance to containment
of the specified isotypic ideal submodules. This closes the precise
higher-jet gate stated conditionally in G-14.

\paragraph{I-27. A Boolean face in the cocharacter lattice.}
The Pl\"ucker ideal of the D-FGA Grassmann cone supplies integral linear
equations in the diagonal cocharacter lattice of D-S3. Under the actual
base-six code embedding, bounded integer deficits cannot carry between
digits. Its intersection with that lattice condition is consequently a
coordinate face of the seed cube. The induced original-bundle characters
then agree on two untested basis directions. The displayed rank-two
unipotent representation detects this equality as exact vanishing of a
nonzero section. No statement about arbitrary decoders follows.

\paragraph{I-28. A hyperplane dimension drop pays for the final parameter.}
The left endpoint defines a curve through its nonzero constant point of
the Grassmannian. Including it in the nested projective-image chain uses
the same dimension budget as the internal bundle automorphisms. If the
constant right functional vanishes on the final image, that image lies
in a proper hyperplane section and loses at least one dimension; this
pays for retaining the right endpoint parameter. D-FGA therefore gives
the sharper $r(w-r)$ support existence statement with both endpoints
included. Its width-three specialization has eight SV pairs, and the
finite flat trace construction of P-12 uses the corresponding seventeen
parameters when sharpened by this result.


\paragraph{R-20. Second-order tangent sections and the exact eventual rectangle range.}
The permanent coefficient morphism has a family of differentials over
the one-parameter motif scheme. Row and column relations cancel its
constant and first-order normal terms. Dividing their polynomial
sections by the square of the parameter gives regular limiting sections
of the ambient coefficient bundle. The finite certificate proves a
nonzero minor after one residue-field specialization; the explicit
degree bounds and cyclotomic argument lift that nonvanishing to the
required characteristic-zero family. D-FGA records the resulting image
dimension, and the D-S3 Fourier characters separate its translated
motifs. The invariant quotient and Veronese construction of G-11 then
give the specified rectangular Schur classes. This transports the
actual proof, including its finite threshold, rather than treating a
dimension inequality as an unspecified positivity assertion.

\paragraph{R-21. Integral image ranks and fixed-rank character extraction.}
The orbit-coefficient homomorphism of free modules has an integral Gram
matrix. Its positive-rank characteristic coefficient is a nonzero
integer, giving the displayed analytic spectral separation. A bilinear
coefficient kernel presents the same image ranks without enumerating
the full tensor basis. The diagonalizable group of D-S3 splits that
kernel into weight blocks; the D-S6 Weyl-character identity extracts
each Schur multiplicity from their ranks. D-FGA identifies the ranks
with the degreewise coordinate-image modules. The algorithm is
polynomial only when the bundle rank is fixed and degree is measured
numerically. The metric estimate is ordinary linear algebra, not an
unproved arithmetic slope transfer.

\paragraph{I-29. A persistent active subbundle and one passive rank.}
On the preserved rank-profile chart, constant transition maps split
into an identity on a rank-$r$ active subbundle and invertible or zero
maps on a rank-$(w-r)$ complement. At the last zero map, the vanishing
higher minors give a row graph over the active summand. The stable
projection argument bounds the support filtration of its coefficient
image by $r(w-r)+w-1$. The Grassmannian dimension argument for rank
preservation adds $r(w-r)$ interpolation pairs. These are explicit
rank-stratum and filtration calculations in D-FGA and D-S3, valid with
arbitrarily many resets; the polynomial pair budget retains its
separate sampling exponent.

\paragraph{I-30. Endpoint sections extend to local automorphisms.}
A nonzero constant endpoint section has a split coordinate functional.
The displayed rank-one update extends the section to a bundle
automorphism on the principal open where its determinant is nonzero.
The determinant is a unit at the center. Thus all generic and constant
interval ranks are preserved after the endpoints are absorbed in the
word. In D-FGA this is an explicit trivialization on the required rank
chart, and in D-S3 it preserves the original coordinate support
filtration. The two width-three rank stages consequently use two and
four pairs, respectively.

\paragraph{I-31. Irreducible matrix cones and deletion of stable blocks.}
The closure of scalar multiples of a polynomial matrix image is an
irreducible affine cone in D-FGA. The rank-$s$ chart identifies the
image of its last block with the graph of a map over the first $s$
columns. Equality of two consecutive projected cones allows deletion:
an invertible right action gives a closed subcone of the same dimension,
which is equal to the irreducible cone. This geometric equality
preserves every homogeneous minor and composes under multiplication
on disjoint parameter schemes. Rank descent and split rectangular
endpoints yield the stated recurrence in the original rank-$w$ bundle.
No exponential compound bundle is introduced, but the recursive block
factor remains and prevents a uniform polynomial sampling conclusion.

\paragraph{P-13. A finite flat algebra encoding any centered generator.}
After clearing Lagrange denominators, distinct base-$B$ exponents embed
the bounded parameter monomials in a single polynomial variable. The
Eisenstein quotient by $T^{B^{2K+1}}-2$ is finite free over
$\mathbb Z$; its generic fiber is a separable field and the displayed
trace-dual basis recovers every coefficient coordinate. The
annihilator argument identifies the resulting image span with the
original coefficient-image module. D-FGA gives the finite algebra and
module interpretation, and D-S3 gives the tensor-coordinate transport.
This is the ordinary algebra trace, distinct from the Frobenius traces
of D-ET, and does not acquire a purity or RH hypothesis by analogy.


\paragraph{G-16. An auxiliary dual line selects a nonrectangular Schur class.}
Form the Segre--Veronese cone of the degree-$n$ coefficient scheme with
the $a$th Veronese of the dual rank-four bundle. Its degree-$d$
coordinate module is exactly the original degree-$d$ module tensored
with $\operatorname{Sym}^{ad}V$. The D-S3 invariant functor and the
D-S6 Schur decomposition select precisely
$(ad+md/4,md/4,md/4,md/4)$. The auxiliary dual line cancels the padding
weight of the explicit Fermat point, so the enhanced cone has a stable
point for every positive padding exponent. D-FGA and the reductive
quotient locators in G-11 give the quotient dimensions; the separate
action and parameter degrees give the effective Hilbert bounds.
Thus the same dimension gap determines an actual nonrectangular
highest weight and a degree progression. The proof does not infer a
highest-weight label from a total dimension inequality alone.


\paragraph{G-17. Partial flags classify the stretched directions.}
For a partition direction $\nu\vdash ns$, subtract its last row to
obtain the polynomial highest weight $\eta$. The orbit of its extremal
line is the partial flag scheme embedded by the exterior-power Schur
construction. Its degree-$d$ coordinate module is the Cartan component
$S_{d\eta}V$, as follows directly by dualizing the orbit span of the
extremal tensor power. Pairing this cone with the $s$th Veronese of the
coefficient scheme makes the D-S3 invariant functor extract exactly
the multiplicity at $d\nu$, in original coordinate degree $sd$.
Positive Fermat coefficients balance the moment map precisely in the
strict region $\nu_4>0$, $\nu_1>(n-m)s$. D-FGA gives the quotient
dimension comparison and its effective Hilbert bounds. The two
excluded faces follow from the earlier ideal containments, so this
is an exact stretched-direction criterion, with its dimension
hypothesis retained. D-S6 records every exterior, symmetric and Schur
operation used in the construction.

\paragraph{I-32. A common-center rank chart with exponential deletion cost.}
Each singular factor is lifted to two univariate bundle maps whose
product depends only on the difference of their coordinates. On the
diagonal center those differences vanish, and every partial interval
is a product of nested projections between invertible boundary maps.
Thus all interval ranks are already correct on that D-FGA chart.
Nevertheless a homogeneous matrix-coordinate section is the product
of all paired differences. A coordinate subspace meets its nonvanishing
locus only if it retains one coordinate from every pair. This gives
the exact exponential threshold for preserving the image cone, while
a dense D-S3 interpolation map keeps the same section nonzero. The
obstruction concerns this coordinate-deletion criterion only.

\paragraph{I-33. Divisors on the one-dimensional torus and compatible weights.}
The prime-weight substitutions are integral cocharacters of the
split torus in D-S3. Pulling back the rational ratio of two products
gives divisors on the multiplicative line. Distinct prime root orbits
can overlap in at most one point, while the inseparable case is
controlled by integer divisor multiplicities. This forces any lost
ratio to be a character, which the coordinate variants separate.
The subset-sum example identifies the exact failure of closure under
products for the raw family. Interpolating its finite evaluation set
by one affine curve repairs scalar product closure because the
coordinate ring of the line is a domain. D-FGA places this curve and
the coefficient-image restriction; no purity theorem is needed.


\paragraph{G-18. Removing scalar orbit periods and comparing every large degree.}
The auxiliary Schur cone and its D-S6 Cartan coordinate modules remain
those of G-17. The smooth Fermat factor has finite projective stabilizer;
varying its padding line in a common stable open avoids every nontrivial
fixed locus. Unique factorization then gives a free projective orbit.
In the D-S3 central group scheme $\mu_4\subset\mathrm{SL}_4$, the
character exponent $ns-|\eta|=4\nu_4$ is trivial. Consequently two
distinct nonzero scalar multiples of the chosen tensor have distinct
closed orbits. Exact reductive invariants separate them, forcing the
degrees of common nonzero invariant sections to generate $\mathbb Z$.
The explicit numerical-semigroup argument gives sections in every
degree at least $\beta^2$. Multiplication by these sections is injective
in both quotient coordinate domains. D-FGA's standard Veronese and
Hilbert-function comparison can therefore be transported to every
integer degree, with the stated effective thresholds and common
asymptotic exponent. The rectangular central congruence is already
encoded by $ns=4\nu_4$. All stability, exact invariant-functor, and
quotient assertions here are in characteristic zero; an integral
model alone is not an assertion of the same behavior at every prime.

\paragraph{I-34. Local relation modules and coordinate cylinders.}
At one factor position, take the kernel of the univariate coefficient
map from the free module on product terms. A rank-one kernel with full
support is a circuit. The D-S3 prime cocharacters and the divisor
argument of I-33 force every canceled multiplicative ratio off this
position to be a character. A free coordinate separates its character
exponents, and the circuit property excludes every proper relation.
Higher local nullity is reduced by coefficient elimination. The result
gives the explicit bound on the number of free coordinate directions;
interpolation replaces the cylinders by integral cocharacters. In
D-FGA terms, the finite family preserves the coefficient-image module
after extension to the working field. The rank-four flattening example
defines a genuine nonzero tensor in the kernel of every one-coordinate
cylinder and identifies the precise failure of that smaller family.

\paragraph{P-14. Rank-one reset maps and scalar bridge sections.}
On a rank-one bundle-map chart, a nonzero pivot writes the map as a
column times a row divided by that pivot. Clearing all pivots splits
the scalar contraction into bridge sections. Invertible interior maps
and absorbed endpoint maps have identity constants over $F(t)$; the
ascending coefficient-row subspaces have at most $w-1$ strict growth
steps. One common interpolation morphism keeps every nonzero bridge
section nonzero, and their product remains nonzero in its coordinate
domain even when adjacent bridges share a reset variable. Constant
linear annihilators give exact coefficient-row preservation. This is
the D-FGA image-module criterion and the D-S3 dense substitution; no
equality of matrix image cones or purity statement is required. The
balanced rational-evaluation form follows from the stated partial
fractions and the same scalar-annihilator argument.


\paragraph{R-22. Cyclic characters and one open differential-rank condition.}
The permanent coefficient morphism is an integral polynomial map of
affine schemes. R-22 works over a cyclotomic extension in characteristic
zero. The D-S3 diagonal torus restricts to $\mu_m$, whose characters
split the differential coefficient space. A single staircase parameter
space supplies both local row and column cofactor relations. Their
highest components define two subspaces with explicitly computed
coprime-factor intersection. The D-S6 Schur determinant is the nonzero
minor of coefficient restriction; its principal geometric specialization
is a Vandermonde determinant with distinct cyclotomic exponents.
The product of finitely many nonzero parameter polynomials defines one
nonempty open subset on which all bidegrees are simultaneously valid.
An integer grid contains a point of that open after the root of unity
is fixed; translated motifs use the prescribed cyclotomic twists.
The final D-FGA differential minor is nonzero in one common deformation,
so its image dimension is at least the stated bound. The proof can be
spread over the localization where these particular denominators and
minors remain nonzero; it does not assert the same rank in every residue
characteristic. Its quotient and highest-weight consequences use the
characteristic-zero dictionaries of G-10 and G-17--G-18.

\paragraph{The G-17 quantitative corollary.}
The rank of a symmetric power and the six-factor Schur-rank polynomial
bound the ambient bundle rank in the D-S6 construction. Together with
the polynomial action degree, these give the binary-length estimate
for the existing first successful quotient degree. This is an effective
bound on the degrees of sections in the D-FGA graded algebra; it is
separate from computing the section or deciding its multiplicity.


\paragraph{R-23. A single graded part of a derived intersection.}
Let $Z=\operatorname{Spec}(R/J_m)$ be the permanent Jacobian scheme,
and let $L=\operatorname{Spec}S$ be a general four-plane in
$\operatorname{Spec}R$. Their D-Tor derived fiber product has
coordinate complex $(R/J_m)\otimes_R^{\mathbf L}S$.
Its first homology in internal degree $m$ is precisely the dimension
defect $\epsilon_m$. The known full-variable linear resolution gives
the preceding map, and the cyclic coefficient calculation proves
that its $2m-2$ old syzygies remain independent while no new constant
relation appears. The degree-$m$ differential then gives
$c_m=4m^2-2m+2-\epsilon_m$. This is an equality of dimensions of
explicit cohomology groups, not an inference from an alternating
$K_0$ class. The low-size positive defects prove failure of the
unqualified restriction argument. The depth bound rules out a full
regular-sequence justification but leaves the single graded
vanishing question open. D-FGA places the associated generic rank
stratum, and D-S6 distinguishes its actual homology from its
Euler characteristic.

\paragraph{P-15. A nonzero section killed by a dense interpolation morphism.}
The displayed quadratic is a rational coefficient tensor with
minimum tensor-train width three: its three balanced flattenings
have ranks $3,4,4$. The upper-triangular bundle-map factorization
keeps every interval on the invertibility open set. Nevertheless
its pullback under the specified centered Lagrange morphism is
zero. The identities in the three linear coordinates $S,U,V$
prove this directly; the harmonic cross-ratio identity supplies
the separately stated rational-evaluation examples. D-FGA records
the nonzero section and its zero pullback, while D-S3 records the
coordinate choices. The example is over $\mathbb Q(\mathrm i)$
for the fourth-root Lagrange nodes, and uses rational abscissas
for the other convention. No scalar change of coordinates is
silently omitted in comparing these two maps.


\paragraph{R-24. The first zero of the graded derived-intersection defect.}
The specified integer Jacobian minor is a section of an exterior
power of the permanent coefficient-map differential. Its nonzero
residue modulo $101$ proves that its principal open locus has
nonempty characteristic-zero fiber. On that locus the image
dimension is $1124$, the exact D-S3 torus quotient bound.
R-23 then identifies $\epsilon_{17}=0$ as vanishing of the
indicated degree of the D-Tor first homology. This is one
explicit generic-fiber result; no complete Tor independence
or assertion in every residue characteristic follows. The
padded coefficient image and resulting highest weights use
the D-FGA and D-S6 dictionaries already proved in G-17--G-18.

\paragraph{R-25. A single family, canceled tangent sections, and induction.}
The bordered matrix is a morphism from an affine parameter scheme
times $\mathbb A^1_\varepsilon$ to the space of forms. Two
coordinated parameter variations cancel their order-one terms
exactly; division by $\varepsilon^3$ on the punctured parameter
line gives fixed Hessian sections in the same differential image.
Their restrictions to the hyperplane $\ell=0$ are the stated
ternary coefficient columns. The intersection of two nonempty
rank opens in the irreducible parameter space permits one matrix
with both required properties. A single minor with nonzero
constant term proves the dimension increment. This transports
the proof to D-FGA rank loci and tangent maps; it does not infer
rank from an alternating $K_0$ class. The finite integer minor
certifies the antecedent at $m=17$ and hence $c_{18}=1262$.
The uniform nonemptiness of those rank opens for all larger
sizes remains the specific geometric question.

\paragraph{I-35. Base loci on a Segre scheme and the seed diagonal.}
The two boundary vectors define a point of
$\mathbb P^{r-1}\times\mathbb P^{r-1}$, embedded as rank-one
functionals. Equality of their annihilator intersections is
equality of base loci of two linear systems on this Segre
scheme. The incidence dimension argument and D-Proper give
an existential subspace with the same base locus and at most
$2r-1$ generators. It need not be realized by an explicit
substitution on the original variables. For the generic matrix,
the Hankel curve realizes such a linear system, although the
matrix value has rank one. This is the known optimal test-space
dimension cited in P-B16, with the new bridge deduction stated
separately.

Independent seed sets give a product of parameter schemes.
At each replacement the other parameters belong to the
coefficient function field, so the extension-field contraction
hypothesis applies. Identifying those parameters afterward is
pullback along a diagonal; a nonzero section can become zero
under that pullback, as the explicit examples prove. Thus
base change of the coefficient field is distinct from seed
identification. Finally the multiplicative lift sends the
rank-$r$ idempotent corner into a one-dimensional endomorphism
corner; its fixed entry readout would have to be a multiplicative
character, contradicted by the two unipotent matrices. D-S3
places these actions, D-S6 places their exterior powers, and
D-FGA places their coefficient-image modules. No new purity
or representation-rank conclusion is inserted by this language.

\paragraph{G-19. A discriminant line and a specified Schur submodule.}
Over the integer coefficient scheme, the homogeneous form and
its partial derivatives define a closed incidence scheme with
projective point coordinates. Its projection is closed by
D-Proper. On the characteristic-zero fiber its reduced
discriminant equation is a section of coefficient degree
$D=k(m-1)^{k-1}$ and variable-change determinant weight
$\delta=m(m-1)^{k-1}$. Benoist's exact degree theorem and
primitive integer construction supply this line; D-S3 records
its character. The highest-weight inclusion into the plethysm
generates one $S_{(\delta^k)}V$ over $\mathbb C$, in the D-S6
Schur-functor language. Its restriction map to one orbit
coordinate ring is zero and to the other injective. This is
a statement about a specified submodule of the coefficient
algebra, not the rank of its entire isotypic multiplicity space.

The Haar identity is its compact complex realization; Gaussian
orthogonality is a positive norm of its integer pullback
coefficients. Neither norm is asserted to arise from a pure
arithmetic sheaf or motive. The explicit cubic Macaulay minor
defines an open localization where the graded Jacobian map
is onto, so the projective gradient-zero scheme is empty.
Its nonzero residue certifies the characteristic-zero smooth
fiber and hence the displayed discriminant nonvanishing.
Positive padding factors the section and forces a singular
intersection, so the discriminant pullback becomes zero.
This same geometric argument proves the stated limitation.
The row-minimality at size three is the D-S3 highest-weight
subalgebra argument applied to the dense D-FGA four-variable
determinant image. Cote~133, scan~028, supplies the Schur
operation viewpoint; scan~019's rank-versus-character warning
remains essential for the specified-copy distinction.


\paragraph{P-16. Reciprocal coordinates, residue functionals, and a necessary center.}
The four disjoint pole divisors lie on the projective line over
the explicitly bounded number field. Their residue functionals
annihilate polynomials of the stated degrees because of the
orders of vanishing at infinity. The displayed rational identity
then makes one nonzero width-three tensor section vanish on
the unbalanced evaluation image. Passing to the coordinate
torus and applying inversion exchanges numerator and denominator
degree. Multiplication by the full coordinate monomial extends
this pullback to the multilinear coefficient tensor; on each
rank-two variable factor it exchanges the two basis vectors.
Every coefficient flattening is consequently changed only by
row and column permutation. This is a direct equality of ranks
of bundle maps in D-FGA, with D-S3 supplying the torus inversion
and diagonal node rescaling; no rank-monotonicity assumption
about polynomial factors is made.

The simple-root discriminant of the node polynomial makes all
$A'(b)$ invertible in the working field, so the Lagrange and
unbalanced rational maps have the stated exact diagonal
identification. Clearing denominators extends the identity
from the dense torus chart to the whole defined evaluation
family. Adding the common center is a different affine
parameter morphism. Its nonvanishing follows from P-14 because
all displayed transition determinants are nonzero polynomials.
Thus the same section is zero on one image and nonzero on the
centered image. D-S6 places the coefficient tensor operations;
it does not identify these two pullback maps merely because
their target bundle ranks agree. The construction uses the
specified algebraic nodes, and supplies no counterexample
for an arbitrary preassigned node set.


\paragraph{I-36. Pairwise dominant pullback and rank-one exterior maps.}
The centered one-pair morphism is dominant on every coordinate
pair: its determinant over the selector function field is the
difference of two Lagrange basis sections and is nonzero.
Hence its pullback is injective on each such bivariate
coordinate algebra. On a rank-one bundle-map chart, a pivot
writes the map as a column and a row divided by a nonzero
univariate section. Clearing the pivots turns an interval
contraction into a product of univariate endpoint sections
and adjacent bivariate sections. Each survives the same
pullback, and their product survives in the parameter domain.

For factors of rank at most $r$, the $r$th exterior map is
zero or rank one. Functoriality of exterior powers, equivalently
Cauchy--Binet, identifies its interval entries with the
original maximal minors. The first paragraph therefore
preserves each nonzero such section. D-S6, V~2.1--2.2.1,
places the exterior-power construction, D-S3 places its
coordinate action, and D-FGA places the resulting determinantal
rank opens. In width three with singular factors this supplies
the rank-two-preserving center. The four residual interpolation
pairs from I-30 preserve the remaining local rank-one tests,
giving five pairs in total. This proof uses actual nonzero
minors throughout; no numerical rank or purity argument is
substituted for the exact pullback assertion. Arbitrary
variable-dependent invertible width-three factors remain
outside this five-pair refinement.


\paragraph{R-26. Characterwise determinantal opens and a linear derived defect.}
The Jacobi parameters form an integral affine scheme; restricting
to nonzero diagonal and edge parameters gives an integral torus
chart. The finite cyclic action on the ternary variables decomposes
the coefficient bundle over characteristic zero into its character
subbundles. In each ordinary character, eight specified normal
sections have a nonzero determinant on a nonempty open chart.
The seventeen finite cores certify those opens by reduction modulo
$101$; their denominators are invertible at the stated certificates.
These are exact algebraic determinants, not approximate ranks.

The cut argument takes place over a Laurent-series field: divide
the second weight component by its path product, replace each cut
edge by a parameter times a chosen nonzero constant, and specialize
the resulting regular normalized matrix at parameter zero. Equal
outer cut orders allow its path-product ratio to have any specified
nonzero limit. Inserting isolated vertices multiplies the indicated
component sections by the common coordinate $x$, an injective map
of the coefficient modules. The two insertion rules therefore
transport each certified determinantal open to all required sizes.
The exceptional character sections are independent by the
coprimeness of consecutive continuants. Irreducibility of the
parameter scheme makes the finite intersection of all these opens
nonempty. This is the scheme-theoretic content of the construction;
it does not assert that one literal zero-edge pencil realizes all
the normalized certificates simultaneously.

Apply the R-25 normal-direction injection and the D-Tor dictionary
of R-23. The dimension of the degree-$m$ first derived tensor
homology is at most $6(m-18)$ for $m\ge18$. Thus the exact remaining
failure of the restricted linear strand is bounded linearly.
D-S3 places the character action, D-FGA places the coefficient
rank opens and orbit-image dimensions, and D-Tor (SGA~6, III~1.5--1.6
and I~5.1--5.4) places the derived restriction. The consequence is a
characteristic-zero bound; the use of modular certificates alone
does not give the whole statement in every positive characteristic.
Combining these dimension bounds with the flag quotients of
G-17/G-18 gives the exact rectangular padding range already for
$m\ge14$, as detailed in R-26. No new positivity conjecture is
substituted for those proved quotient comparisons.

\paragraph{I-37. Rank-two active bundles and the second coefficient form.}
The rational centered two-pole family defines a morphism from an
open parameter scheme to the coordinate affine space. Its one-pole
subfamily is dominant on each coordinate pair. An irreducible
multilinear section can be considered separately because its
irreducible factors use disjoint coordinate sets; specializing the
other coordinates preserves the affine singular factorization.
On the nonzero minor charts, every strict internal product has
rank two. A rank-one internal product would have a polynomial
column--row factorization over the coordinate UFD, and would split
the irreducible scalar into two nonconstant factors.

Exterior squares and the pairwise dominant pullback preserve these
rank-two charts, by I-36. Compatible bases split each intermediate
rank-three bundle into a persistent active rank-two summand and
a passive line, with constant transition $\operatorname{diag}(1,1,0)$.
The affine singular condition makes the linear passive self-entry
zero. If the original section vanished on the two-pole image,
simple-pole residues would annihilate all first coefficient forms
at the one-pole center. Translation along its nonzero common
slope vector annihilates the contraction of the second coefficient
form with that vector. In the compatible bases this form has only
adjacent off-diagonal entries. Induction along the ordered path
then makes each of them zero in every characteristic. Every
possible active-to-passive-to-active output transition consequently
has a zero product, forcing the whole section to be zero.

This restates the complete contradiction in bundle-map language.
The second coefficient form is extracted formally from multilinear
polynomials, so division by two is never used. D-S6, V~2.1--2.2.1,
places the exterior-square maps; D-S3 places the changes of frame;
D-FGA places the rank strata and nonvanishing charts. The final
diagonal node rescaling identifies the rational family with
ordinary centered Lagrange interpolation. Its five parameter
coordinates yield the exact grid count in I-37. Invertible
variable-dependent transitions can propagate the passive line,
so the tridiagonal argument does not extend to that unresolved
mixed-rank case.


\paragraph{I-38. Triple-dominant pullback and isolated invertible transitions.}
The one-pole parameter morphism is dominant on every three
coordinate factors. Indeed, after solving for the center and
amplitude using two coordinates, the third is a nonconstant
rational function of the selector over their function field:
three distinct simple poles forbid both interpolation coefficients
from being constant. This argument establishes dominance even
when a characteristic-two Jacobian test is inseparable.

For the promised alternating rank pattern, exterior-square maps
of the rank-two factors have rank one, and there is at most one
invertible exterior-square map between two of them. Each scalar
connector is therefore a section depending on at most three
coordinates. Triple dominance preserves every required nonzero
connector. Compatible active rank-two summands and passive lines
then put the constant maps in the forms $I_3$ and
$\operatorname{diag}(I_2,0)$. The rank pattern survives the
irreducible multilinear reduction: a specialized singular factor
between two surviving invertible ones is absorbed into the latter
and makes it singular.

The second coefficient form at the one-pole center annihilates
both vectors $h_i=(a_i-y)^{-1}$ and $h_i^2$. Its only possible
entries are within two positions of the diagonal. The determinant
$h_jh_k(h_k-h_j)$ is a nonzero section, in every characteristic,
so row induction annihilates the whole form. Every passive
excursion of an output path has one of its zero entries as a
factor. This transports the contradiction of I-38 without
asserting preservation of each individual minor, and without
discarding passive propagation at invertible factors.
D-S6, V~2.1--2.2.1, supplies exterior powers; D-S3 supplies the
bundle changes of frame; D-FGA supplies the rank-open and
function-field formulation. Longer consecutive invertible runs
require connectors on four or more coordinates and remain outside
this triple-dominance proof.


\paragraph{R-27. Formal endpoint jets and exact derived restriction.}
The endpoint deformation gives a matrix of coefficient sections
over the formal parameter ring $\mathbb C[[\varepsilon]]$.
Its first and second normal maps are the successive leading
coefficient maps after eliminating the constant-rank image.
The lower triangular coefficient matrices represent its
reductions modulo powers of $\varepsilon$. Their ranks therefore
count the unit, first-order and second-order pivots exactly.
The boundary projection retains complete character subbundles.
The support of the two successive normal maps lies wholly
inside those subbundles, so the projected quotients agree
with the full quotients; discarded character sections cannot
create an additional correction.

The fixed size-seventeen certificate proves nonempty opens
for the three coefficient ranks. Central insertion transports
its normalized sections by the injective monomial maps
$f\mapsto xf$, $f\mapsto yf$, $f\mapsto zf$ on disjoint
weight windows. The cut is interpreted in a Laurent-series
chart after canceling the long-path edge product. The exact
endpoint identities give a separate upper bound of four
on the first normal rank. This upper bound is required:
lower rank bounds alone would not control the second normal
map under specialization. Intersecting the rank opens and
this identity chart gives all six missing normal directions
at a single actual pencil. R-25 then makes the restriction
complex exact in the specified degree in every size at least
seventeen. D-S3 places the character action and changes of
frame, D-FGA places the determinantal opens, and D-Tor
(SGA~6, III~1.5--1.6 and I~5.1--5.4) places this precise
derived-tensor vanishing. No regular-sequence assertion in
all degrees is made.

\paragraph{G-20. Exact comparison of quotient Hilbert dimensions.}
Insert R-27's exact image dimension into the enhanced
partial-flag quotients of G-17/G-18. Their dimension difference
is the integer $C_m-2n^2+1$. The two standard-Veronese
Hilbert polynomials have positive leading coefficients;
their original-degree ratio therefore has exactly that
growth exponent. The exponent is never zero for the integer
pairs in the theorem. This removes the former unresolved
dimension transition, without determining every finite
graded piece. D-FGA and EGA~II, 2.5.2, place the projective
coordinate rings; the invariant-quotient and stabilizer
locators are the same as those already recorded for G-17/G-18.
The first-row condition for rectangular weights converts
to $n<4m/3$. The size-thirteen example proves sharpness of
the threshold fourteen for the uniform eventual statement,+not absence of an individual earlier-degree equation.

\paragraph{G-21. Additional-variable tangent summands.}
The cofactor sections form an independent rank-$m^2$
subbundle on a nonempty quaternary parameter open. Intersect
it with the exact differential-rank open of R-27. Extending
the coefficient bundle by unused variables adds one copy
of that cofactor subbundle in each new variable multidegree.
These summands have zero intersections with one another
and with the old tangent image. The resulting rank is
$rm^2-2m+2$, equal to the row--column orbit-parameter bound.
The same degree-$m$ tensor complex from D-Tor has kernel
of dimension $2m-2$, exhausted by the old linear relations.
Independence of the cofactors handles degree $m-1$; lower
degrees have zero source. This proves the exact graded
homology statement, including its stated limits. For $r>m^2$
the coefficient specialization is an algebra morphism into
a larger polynomial ring, not a quotient by a negative
number of equations.

\paragraph{I-39. Failure of rank-stratum preservation and nonintegrable jets.}
The ordinary centered one-pair morphism sends the explicitly
given generic rank-two interval to the rank-one stratum,
even though the scalar coordinate section is irreducible
and stays nonzero. The two singular anchors and the unipotent
interior give an exact determinantal equation for this
specialization. Thus preservation of every internal rank
stratum is a stronger invariant than scalar nonvanishing.
D-FGA supplies the determinantal-stratum language, while
D-S3 places the affine unipotent transition factors.

The separate Hessian family is a nonzero symmetric section
with the asserted tangent kernels and rational-curve
quadratic restriction. It cannot be the second coefficient
pullback of a fixed polynomial along the whole one-pole
parameter morphism: restriction to zero amplitude would
force independence of the selector, contrary to its explicit
entries. This is an exact failure of the proposed abstract
jet criterion. It is not a counterexample to the two-pair
scalar nonvanishing problem; the remaining construction
must retain the coordinate-ring and integrability conditions.


\paragraph{R-28 and G-22. Added tangent bundles and Fourier-frame stability.}
The determinant differential has an independent cofactor subbundle
of rank $n^2$ on a nonempty quaternary parameter open. Each added
variable contributes a copy in a separate multidegree. The resulting
rank $(r-2)n^2+2$ equals the parameter dimension minus the dimension
of the generically finite-stabilizer left--right special-linear
orbit. The permanent padding map is generically injective after
projectivization: the irreducible nonlinear factor and padding
line are recovered uniquely. D-S3 identifies the group actions,
and D-FGA and EGA~IV$_3$, 13.1.3, identify these rank opens and
fiber-dimension calculations.

In G-22, the partial Fourier matrix is a frame in a rank-$r$
locally free representation, after scalar extension to $\mathbb C$.
Its rank-one Hermitian moments sum to the required diagonal
weight. The ordered tensor and the extremal flag line have
zero total traceless moment and finite projective stabilizer.
Multiplication of the linear sections and the outer symmetric
power preserve the minimum weight for every cocharacter, since
the corresponding initial sections have nonzero products in a
symmetric algebra. The Hilbert--Mumford criterion therefore
transports stability to the multiplied section. This is a
cocharacter and initial-term dictionary in D-S3; it does not
assert equality of the two Hermitian moment maps. The Schur
and exterior tensor operations belong to D-S6, V~2.1--2.2.1,
V~3.2 and V~3.10, with the precise representation-ring convention
already fixed above. The square roots give a complex existence
witness, not an integral point or an effective arithmetic sampler.

The cyclic product-plus-power family is an irreducible parameter
scheme. Nonzero coefficient minors impose the finite projective
Lie-stabilizer condition. Irreducibility of a form is open because
the reducible locus is a finite union of proper projective
multiplication images. An explicit Eisenstein section and the
coefficient calculation of its stabilizer show these opens are
nonempty. Their intersection with the stable open remains nonempty.
Varying the padding line avoids the fixed loci of the finite
automorphism group and leaves only the central group scheme
$\mu_r$ upstairs. Its character on the enhanced representation
is $r\nu_r$ and is trivial. Thus distinct nonzero scalar multiples
have distinct closed orbits. Invariant separation makes the
nonzero evaluation degrees have gcd one; the numerical-semigroup
argument gives every degree from $\beta^2$. Proper projective
images use D-Proper, and the character and stabilizer statements
use D-S3. All these orbit assertions are over characteristic zero.

\paragraph{G-23. Exact growth of enhanced invariant coordinate algebras.}
The Cartan cone attached to the flag weight has degree-$j$
coordinate module $S_{j\eta}V$. The diagonal Segre--Veronese
algebra is therefore
$\mathbb C[Y]_{sj}\otimes S_{j\eta}V$; taking invariants
extracts the single label $j\nu$. The stable quotient dimensions
are $u=(r-2)n^2+r_\eta-r^2+3$ and
$v=rm^2-2m+r+r_\eta-r^2+2$. This is a dimension computation
on two projective invariant schemes, not an equality of their
coordinate modules in $K_0$.

The explicit standard Veronese has $E=rT\beta!$.
Its positive leading Hilbert coefficients, combined with
nonzero sections in all sufficiently large original degrees,
give the ratio $\kappa j^{v-u}(1+O(j^{-1}))$ through every
integer stretch. The character-ring reduction cannot remove
the difference $v-u$; its nonvanishing is a separate elementary
congruence argument. D-FGA and EGA~II, 2.5.2, place the graded
coordinate algebras and projectivizations. The coefficient-degree
bound on the action and the parameter-degree bound on the orbit
map remain distinct. The construction proves the eventual sign
in the stated strict cone with $4\le r\le m$; boundary faces,
more rows, and effective positive extraction at individual
finite degrees remain outside this conclusion.

\paragraph{P-17. Exterior decoders, coefficient-rank strata, and odd weights.}
On the integral node scheme obtained by inverting all node
differences, the three-row Vandermonde matrix is a split
surjection and its kernel $K$ is locally free of rank $n-3$.
The cubic coefficient decoder is a diagonal isomorphism with
unit Vandermonde factors. Its balanced-ideal part is
$\bigwedge^3 K$, by the explicitly cited ideal theorem.
The coefficient-cut rank bounds are determinantal closed
conditions. The three-block contradiction, the support-at-least-four
property, and the leading exterior projections are valid in
every residue characteristic on this node scheme. The two end
kernel subbundles have intersection ranks zero and one at
$n=9$ and $n=10$, yielding respectively the incompatible wedge
degrees and the forced common exterior divisor in the proof.

For an invariant multilinear tensor, each Schmidt support is
an $\mathrm{SL}_2$ subrepresentation on its fixed-rank chart.
At an odd cut its split-torus characters are nonzero odd
integers. The Weyl element pairs opposite character spaces, so their total rank is even, including over characteristic two.
With rank at most three this forces rank two at every odd cut,
and each affine core factors through a rank-two bundle.
The two-pair theorem I-37 then applies. Translation and reversal
are the two unipotent operations on the multihomogenization;
the closed coefficient-rank degeneration to the leading
translation term gives the eight-variable closure. The exact
torus reference is SGA~3, I~4.7.3; exterior and symmetric
operations use D-S6, and the rank-stratum formulation uses
D-FGA. No division by two or semisimplicity in positive
characteristic is used.

\paragraph{I-40. Determinant characters of Borel-stable support bundles.}
For the homogeneous common-translation-invariant tensor, the
two support bundles at a cut are Borel-stable, although they
need not be $\mathrm{SL}_2$-stable. Their determinant lines
are upper-unipotent fixed lines inside exterior powers of
full rational $\mathrm{SL}_2$ representations. The standard
dominant-character fact cited from Jantzen, Lecture~I,
pp.~2--3, gives nonnegative integer weights $\delta_i,\delta_i'$.
The nonsingular coefficient pairing of total weight $m=n-2D$
gives $\delta_i+\delta_i'=r_i m$. This is an equality in the
character lattice of the split torus, not a congruence in
the residue field.

The support inclusions $L_i\to L_{i-1}\otimes V$ give
equivariant affine core maps. Their unipotent covariance
makes every nonzero square determinant constant in the local
variable. Its constant term is then invertible, and its
torus covariance gives the exact increment
$\delta_i-\delta_{i-1}=r_i$ across a full-rank square core.
Telescoping the determinant-line character proves the
full-width run bound $n-2D$. Constant absorption preserves
it because a removed singular map cannot leave the next
full-width map invertible. At width three, weights zero and
one put the word in I-37 and I-38 respectively. The cubic
closure of P-17 then makes a nine-variable quartic annihilator
translation invariant of weight one and gives the contradiction.

This restatement uses the determinant exterior operation in
D-S6, the split torus and Borel group schemes in D-S3, and
the coefficient-rank charts in D-FGA. The rational evaluation
map is regular after inverting the node-pole differences.
Its generic two-simple-pole chart is dense in the balanced
degree-two rational family even in characteristic two.
Clearing the squarefree coordinate denominator and using the
explicit selector and Newton-simplex grids gives the finite
sample count in I-40, with its working-field and extension-input
conditions retained.


\paragraph{G-24. Boundary cocharacters and restriction of invariant algebras.}
At $n=ra$ and $m=(r-1)a$, the cocharacter with weights
$(r-1,-1,\ldots,-1)$ sends a padded form to
$x_0^a q(0,x_1,\ldots,x_{r-1})$. Every invariant section
has the same value on this limit. Restriction therefore
injects the padded invariant coordinate algebra, with its
original grading, into the $(r-1)$-variable permanent
invariant algebra. Injectivity follows from the limit on
the dense padded image, not from an assumed equality of
the two algebras. This is a cocharacter-specialization
statement in D-S3 and a morphism of graded coordinate
algebras in D-FGA.

For $r\ge5$ the lower algebra has the exact stable quotient
dimension supplied by G-21. For $r=4$ it is bounded degree
by degree by the full ternary invariant algebra, whose map
onto the lower permanent algebra is a surjection by
exactness of invariants. R-28 supplies a separate free
stable point in each comparison ring by adding one linear
factor; the integer $a\ge2$ makes both auxiliary directions
strict. Thus both degree semigroups have gcd one. Their
dimensions satisfy $u>v$, and the explicit common generator
bound and Veronese make the Hilbert comparison strict at
every sufficiently large original degree. Above the
boundary the same cocharacter sends the form to zero.
Together with G-23 this proves the exact eventual rectangular
range through $m$ rows, with the stated finite-degree and
nonrectangular exclusions from its scope. EGA~II, 2.5.2,
places the projectivizations, and the stable quotient and
image-degree locators are those already recorded for G-23.


\paragraph{R-29 and G-25. Two-product families and enlarged stable quotient schemes.}
Work over $\mathbb C$ with a rank-$r$ vector space split as two
summands of ranks at most $m-1$. The ordered linear-factor tensors,
their symmetric products, and the auxiliary Cartan flag cone are
formed by the tensor and Schur operations of D-S6. Minimum
cocharacter weights are preserved under multiplication in the
symmetric algebra. Kempf--Ness and unitary transitivity on the
complex partial flag variety give a product section with any
prescribed positive rational moment spectrum in each summand.
This analytic normalization is an existence argument on the
complex fiber; it does not assert an integral section or an
effective arithmetic change of basis.

The orthogonal sum of the two product sections has the required
moment spectrum. The Lie-stabilizer equations split by bidegree:
mixed blocks vanish by independence of directional derivatives,
and diagonal blocks are scalar. The affine scheme parametrizing
two products and a padding line contains the stable open, the
irreducible-form open, and the finite-projective-stabilizer open.
Their nonempty intersection and avoidance of finitely many fixed
line loci give a common stable section with only the central
projective stabilizer. Its central character is $r\nu_r$ and
is trivial on $\mu_r$. These are group-scheme and character-lattice
statements in D-S3 on the characteristic-zero fiber; proper
projective multiplication images use D-Proper.

The degree-$j$ Cartan coefficient algebra still extracts precisely
the label $j\nu$. Hence the projective invariant schemes have
the same dimensions $u,v$ as G-23 for $r\le2m-2$. G-25 checks
the only changed dimension lower bounds and carries across the
same standard Veronese, positive Hilbert leading terms and
gcd-one degree semigroup. The exact difference is unchanged.
For the new positive family $r=n=m+1$ it equals one, and the
label is $(3j,(2j)^{m-1},j)$ in degree $2j$.
D-FGA and EGA~II, 2.5.2, place these graded algebras and their
projectivizations. This comparison of actual invariant Hilbert
functions is stronger than exhibiting one equation copy, but
does not provide a general positive formula for individual
finite-degree multiplicities.

\paragraph{I-41. Saturated weights, divided squares, and pole ranks.}
On the fixed-rank coefficient charts of I-40, a width-three run
of two invertible cores exhausts the determinant-character budget
at weight two. The bounding determinant lines have character zero
and are fixed by the whole split $\mathrm{SL}_2$. Their decomposable
exterior tensors determine invariant rank-three support bundles.
Single-coordinate contractions give neighboring invariant
rank-two bundles. Torus character parity and the Weyl involution
force weights $1,-1$: larger opposite weights would leave no
odd-rank invariant subbundle in their tensor product with the
standard bundle. This uses the integral character lattice of
SGA~3, I~4.7.3, rather than reduction of weights modulo the
residue characteristic.

The invariant rank-three subbundle in $V\otimes V$ is explicitly
spanned by $e_0e_0$, $e_0e_1+e_1e_0$, and $e_1e_1$.
It is the divided-square tensor submodule in every characteristic.
The proof retains its actual dual coefficient pairing; it neither
identifies it with its dual in characteristic two nor assumes
semisimplicity. Determinant and tensor operations use D-S6,
and the support charts use D-FGA. The character shifts make
the two middle affine maps upper unitriangular in compatible
weight bases. Their two-dimensional connector determinant has
the nonzero square coefficient $d_2d_3$.

Over the node scheme with all node differences inverted, the
one-pole rational evaluation is regular after also inverting
the node-pole differences. The unique double pole of that square
term proves nonvanishing of each saturated connector. Compound
rank-one factorization then preserves every required interval
rank. On the resulting compatible charts the second-coefficient
matrix has bandwidth three and annihilates the three sections
$1,h,h^2$. Their three-column Vandermonde determinants are units,
so the row induction and passive-path argument force the
contradiction in every characteristic. Translation and reversal
are root-group operations; their leading limits stay in the
closed determinantal width locus. They reduce the ten-variable
degrees to the weight-two and weight-zero cases. The balanced
finite interpolation grid is the same regular-open construction
as in I-40, with all field-size and extension conventions retained.


\paragraph{R-30 and the G-25 endpoint extension. Compact forms and auxiliary flags.}
At a block of rank $m$, the nonconstant torus character defines
a nontrivial partial flag. Its stabilizer intersects the Fourier
basis torus only in scalars: the corresponding circulant matrix
cannot preserve a proper consecutive coordinate subspace unless
every nonconstant cyclic coefficient vanishes. This produces a
stable ordered tensor even though its polynomial product alone
has a positive-dimensional complex stabilizer. After normalization,
the product has only a finite compact projective stabilizer.
The tilted padding line subtracts a positive rank-one Hermitian
operator from the desired moment matrix, with two explicitly
chosen rational positive eigenvalues. Both product blocks then
have nonconstant spectrum.

The full ordered padded tensor is the point at which the moment
vanishes. Losev's real-form stabilizer calculation applies there
and makes its complex stabilizer finite. The two-product parameter
scheme also has a nonempty open with finite padded-form stabilizer.
A generic section of the nontrivial partial flag orbit avoids
every fixed locus of that finite projective group. The resulting
enhanced point has no scalar period. This is a characteristic-zero
stable-quotient construction using D-S3 group stabilizers and
cocharacters, D-S6 tensor/Schur operations and D-FGA flag and
coordinate schemes; the Hermitian normalization is not transported
as an integral arithmetic algorithm. The exact G-25 quotient
comparison consequently extends to $r\le2m$ with unchanged
explicit constants.

\paragraph{G-26. A Pieri face of the bidegree pullback.}
The padding parameter morphism pulls a degree-$d$ coordinate
section back to bidegree $(ad,d)$ on the line and the permanent
cone. Its injectivity comes from dense image. In the representation
Grothendieck group its Schur components are governed by
$[\operatorname{Sym}^{ad}V^*]\,[S_\alpha V^*]$ and the
Pieri identity. The direct-sum multiplicity comparison is taken
over $\mathbb C$, where these representations are semisimple;
no integral direct-sum decomposition is inferred from a $K_0$
identity. When $\lambda_1=ad$, the horizontal strip must use
every column, so the unique surviving label is the tail of
$\lambda$. The zero-last-character coefficient subalgebra
identifies that tail multiplicity in one fewer variable.
These tensor and Schur identities belong to D-S6, with the exact
Pieri source locator retained in G-B20.

Two different enhanced invariant schemes now provide the
comparison: the determinant scheme in $r$ variables and the
tail permanent scheme in $r-1$ variables, or its full ternary
ambient bound at $r=4$. Their separate stable points have trivial
scalar period. Deleting the first flag part removes a nonnegative
number of strict weight comparisons, and the exact quotient
dimensions satisfy $u>v$ throughout the equality face with
$4\le r\le m$. The standard Veronese and the two separate
degree semigroups give strict comparison in every sufficiently
large original degree. D-FGA and EGA~II, 2.5.2, locate these
graded coordinate algebras and projectivizations. The argument
uses an injection and an upper bound, not a surjective lifting
of lower-dimensional sections to padded forms.

\paragraph{P-18. Three block characters and adjacent contraction degrees.}
The blockwise homogeneous coefficient tensor is graded by
three independent split tori. The three summands of its common
translation derivative have different block characters, hence
each is zero in every residue characteristic. At an extreme
block cut, the last-variable coefficient and the constant-prefix
coefficient occupy different degree blocks; their ranks add.
This forces two-dimensional extreme support bundles. The
prime rational-evaluation kernel makes every nonzero middle
coefficient map between those bundles invertible, and an
interior cut has two disjoint rank-two blocks. It cannot lie
in the width-three determinantal locus.

For a smallest unrestricted cubic, the exterior decoder of P-17
lives in $\bigwedge^3K$ over the node scheme. Decomposition by
left exterior degree makes the three contraction operators
connect only adjacent degree blocks. A zero block therefore
separates a smaller kernel element. Minimality and the P-17
rank lemma leave only adjacent occupied degrees at every
prefix; the ordered support is consequently trilinear on
three intervals with at most two shared boundary variables.
The degree-three interpolation grid has $20(n+1)^2$ points.
This dictionary uses the torus characters in D-S3, exterior
operations in D-S6 and determinantal charts in D-FGA. It neither
changes the variable order by a general linear coordinate map
nor asserts that setting a shared boundary variable to zero
preserves the balanced ideal.

\paragraph{I-42. Rank-one reset blocks and the special fiber at two.}
On the compatible rank charts of I-41, the mixed second
coefficients factor as $a_i b_j$ inside each reset block and
vanish across a strict internal reset. These are polynomial
rank-one equations; they retain more information than a
bandwidth condition alone. The three tangent sections
$1,h,h^2$ give a split rank-three evaluation map on any three
distinct nodes. The first nonzero row and the first future
nonzero coefficient reduce a hypothetical relation to at most
three Vandermonde columns when the reset distance is at most
five. At distance six the remaining relation yields exactly
the factor $2a_kb_\ell$. This also vanishes only trivially on
the open base where two is invertible.

The characteristic-two fiber is genuinely different: two
four-column kernels overlapping in one node give the sharp
seven-node abstract counterexample. Disjoint four-column
kernels give the eight-node example in every characteristic.
Neither is asserted to be the jet of an actual annihilating
ROABP. Use mixed coefficients, rather than division by two,
for this second-order statement. The fixed-rank and kernel
schemes belong to D-FGA, while the split characters and
divided-square anchors of saturated runs use D-S3 and D-S6
as in I-40/I-41. A saturated run exhausts the determinant
character budget, so the same anchor square coefficient
preserves its one-pole rank regardless of its length. The
sharp reset bounds then prove the stated irreducible classes
at weights three and four, and weight five away from two;
shorter nonsaturated runs retain their explicit rank gate.


\paragraph{P-19. A three-coordinate transversal and an injective exterior projection.}
Over the distinct-node base, the rank-three Vandermonde map
defines a kernel bundle $K$. Every coordinate subbundle on at
most three nodes meets $K$ trivially; projection away from such
nodes is consequently injective on $K$ and on its third exterior
power. The P-18 ordered-support reduction and two disjoint
rank-two coefficient blocks show that a smallest width-three
cubic counterexample would have all its monomial supports
meeting one set of at most three coordinates. Its decoded
section of $\bigwedge^3 K$ would be killed by that injective
projection. This contradiction works on every field fiber of
the distinct-node base. The proof uses determinantal coefficient
charts from D-FGA, split-torus degree blocks from D-S3, and the
exterior operations of D-S6. It closes both boundary cases of
P-18, with $20(n+1)^2$ evaluation points, and leaves degrees
four through $n-4$ as the remaining balanced width-three gate.

\paragraph{R-31. A residual scalar character on a common parameter scheme.}
On the open parameter scheme where the $2m+1$ product and
padding factors are independent, a compensated transposition
fixes the polynomial and its auxiliary covector. Its determinant
character is $-1$. After lifting to the special linear group,
the enhanced tensor has character $(-1)^b$. Density forces
every odd graded invariant to vanish when $b$ is odd. The
coordinate ring and its Cartan extraction are the D-FGA and
D-S6 constructions of G-22; the stabilizer character belongs
to D-S3. Over $\mathbb C$, closed-orbit separation then rules
out a nonzero closed point with trivial scalar period on this
common enhanced subcone. It does not impose the same character
on the larger permanent or determinant schemes. The obstruction
therefore narrows the next stable-point construction rather
than supplying an equation or multiplicity comparison for the
full cones.


\paragraph{R-32 and G-27. Unpadded stable quotients and a one-sided degree semigroup.}
The same two-product parameter morphism, without a padding
line, meets every prescribed positive moment direction when
each factor block has rank at most $d-1$. Its enhanced tensor
is stable. A nonzero flag weight allows a generic flag section
to remove the remaining finite projective stabilizer. The
central character then vanishes on the special linear center.
For a rectangular label the flag is a point, so only stability
and the quotient dimension are retained. This is the D-S3
stabilizer and cocharacter construction with the D-S6 Schur
line and D-FGA enhanced coordinate scheme; no integral
moment-map assertion is made.

G-27 uses two such schemes, in ranks $r$ and $r-1$. The upper
comparison needs no degree-period-one property on the lower
scheme. If $g$ is the gcd of its homogeneous generator
degrees, every sufficiently large multiple of $g$ is present.
Whenever its degree-$j$ piece is nonzero, a nonzero multiplier
places that piece inside a suitable degree of the standard
Veronese. The determinant algebra, used for a lower bound,
does have period one. Thus the projective Hilbert comparison
works at every original degree without identifying the two
semigroups. These are the graded-domain and Proj constructions
of D-FGA and EGA~II, 2.5.2. Together with the Pieri tail
dictionary of G-26, they classify every fixed nonrectangular
direction through $2m-1$ rows. The finite degree thresholds
are effective bounds, not polynomial-size algorithms.

\paragraph{I-43. A Pl\"ucker line constrains actual adjacent support bundles.}
A determinant-weight-one line in the third exterior power
generates a standard rank-two $\mathrm{SL}_2$ module. Its
projectivization lies in the Grassmannian, whose line
description gives an invariant rank-two subspace inside an
invariant rank-four space. The exact known Grassmann input
is I-B10, Lemma~3, p.~4. The actual preceding coefficient
support has rank at most three. Its split-torus characters
and the physical unipotent action exclude a larger pair
of opposite weights in that invariant two-plane, leaving
weights $(1,1,-1)$ on the rank-three support.

No splitting of these module extensions is needed in a
positive-characteristic fiber. The normalized core updates
have a common covector, so their second exterior map is
affine in the internal variables. A single node pole detects
each nonzero coefficient while the two anchor evaluations
remain independent. The determinant-weight-zero case uses
one divided-square anchor and an arbitrary opposite affine
anchor. These identities hold on the actual coefficient
rank charts over the same rational parameter field. The
preserved two-dimensional contractions compose, including
the endpoint projections, before applying the I-42 reset
argument. D-S3 locates the group scheme and characters,
D-S6 the exterior/divided operations, and D-FGA the
Grassmann and determinantal charts. The result is valid on
every field fiber with distinct nodes: all weights at most
three and, using P-19, every multilinear input on at most
eleven variables are hit by the balanced two-pair sampler.


\paragraph{R-33 and G-28. A generic Pl\"ucker chart and a quotient upper bound.}
The maximal differential rank of the permanent coefficient
morphism identifies its generic tangent kernel with the
row-column diagonal gauge directions. On the Grassmann
open where every Pl\"ucker coordinate of the coefficient
matrix span is nonzero, any diagonal operator preserving
that span has equal sums on every fixed-size coordinate
subset. Exchanging one index forces all individual weights
equal. The zero gauge character makes the scalar zero.
Thus the generic projective stabilizer has zero Lie algebra
and is finite over $\mathbb C$.

These are the determinantal and Grassmann opens of D-FGA
and the torus character calculation of D-S3. Stabilizer
openness retains the precise SGA~3 VI$_B$, Proposition~4.1
and cited EGA~IV$_3$, Theorem~13.1.3 framework. The generic
fiber dimension equality is separately cited from the
Stacks Project, Lemma~33.20.4(2), tag~0B2L (R-B18), whose
complete primary proof was read. A fiber of the invariant
quotient contains the orbit through its generic point.
This gives only an upper bound on quotient dimension;
it asserts neither stability nor equality.

G-28 uses exactly that upper bound on the lower permanent
algebra. Its actual dimension may be smaller, or the ring
may contain only constants. The projective image degree and
Hilbert upper estimates decrease with that dimension, so
the G-27 comparison still holds. This closes all fixed
nonrectangular directions through $2m$ rows and extends
the boundary-only comparison to
$r\le\min\{2n-2,m^2\}$. The Schur tail, Proj and degree
semigroup constructions retain the D-S6/D-FGA interpretation
already given; generic freeness is not substituted for
the stable-point hypothesis needed for the upper determinant
algebra's period one.

\paragraph{I-44. Torus chains, anchor charts, and a residue-support condition.}
The normalized cores move between integer torus characters
differing by two. Unless the rank-three support is a chain
of three successive such characters, their updates have
a common image or covector, and the second exterior map
is affine. In the chain case, the actual preceding support
allows four weights. Minimal coefficient independence and
rank two leave precisely the five anchor charts displayed
in I-44. These statements use the D-S3 character grading
and D-FGA rank charts and hold on every characteristic fiber.

Two anchor types have a unique nonzero pole coefficient.
The third, paired with the actual reverse anchor, reduces
the second exterior map to an explicit quadratic. A square
coefficient gives a double pole. In the multilinear case,+the coefficient of $tz$ and the simple residues impose two
independent equations on each row. Every nonzero row would
therefore need at least three neighbors. With at most five
internal positions, one of the two coefficient types has
at most two entries, forcing an order that makes the
quadratic factor into nonzero linear forms. The evaluation
kernel cannot contain that product.

The resulting interval and endpoint assertions use a single
parameter field before invoking the reset-Hessian theorem.
The exterior operations belong to D-S6, without division
by two. The characteristic-two distinction occurs only in
the subsequent I-42 scalar Hessian step: four-core runs
are covered on every fiber, and five-core runs where two
is invertible. Factor weights then give the complete balanced
theorem through twelve variables in every characteristic
and thirteen away from two, with the stated finite grid.

\subsection{What the manuscript record does and does not establish}
The supplied cote~133 study also reports decisions about the name on scan~008, a characteristic-two pre-$\lambda$ issue on scan~006, and a crossed-out Stiefel--Whitney assertion on scan~009. Those reports are retained as reports of earlier readers, with their own scope. No new visual decision about the name or the deleted text is made in this run. The five attached LaTeX transcriptions contain no original raster pages. The separate transcription table records the forced algebraic readings and the distinction between a textual ambiguity, a mathematical correction, and a claim whose ink cannot be checked here.
